\documentclass{report}
\usepackage[utf8]{inputenc}
\usepackage[french]{babel}
\usepackage{amsmath, amssymb, amsthm}
\usepackage{tikz}
\usepackage{geometry}
\usepackage{hyperref}
\geometry{margin=1in}
\title{Jalon 10 : Changements de base}
\author{Charles EDOU NZE}
\date{\today}
\begin{document}
\maketitle
\tableofcontents
\chapter{Cours : Changements de base}
\subsection*{1. Présentation du concept clé}

Imaginez un cartographe qui, pour la première fois, explore une région inconnue. Il pourrait commencer par établir un système de coordonnées simple, basé sur les points cardinaux : "tant de kilomètres au nord, tant à l'est". C'est une base de référence naturelle, intuitive. Cependant, si cette région est dominée par une chaîne de montagnes orientée obliquement par rapport à son axe nord-sud, ou traversée par un fleuve sinueux, décrire les caractéristiques géographiques importantes (les crêtes, les vallées, le cours du fleuve) dans ce système initial peut s'avérer laborieux. Les coordonnées seraient complexes, les équations des lignes de force de la topographie alambiquées.

Le cartographe avisé réaliserait bientôt qu'il serait bien plus efficace d'aligner son système de coordonnées principal le long de la chaîne de montagnes, ou de l'axe général du fleuve. Les descriptions des éléments clés deviendraient alors d'une simplicité déconcertante : "le sommet est à telle altitude le long de l'axe de la crête", "le fleuve suit l'axe principal avec de légères déviations". Le paysage n'a pas changé, mais la *manière* de le décrire a été transformée. Ce processus de réorientation du cadre de référence est l'essence du \textbf{changement de base}. La \textbf{matrice de passage} est l'instrument précis qui permet de traduire les coordonnées d'un point du système initial vers le nouveau système, et inversement. C'est un dictionnaire universel, une clé de voûte pour passer d'une langue descriptive à une autre, sans altérer la réalité sous-jacente.

Dans le même esprit, considérons une carte topographique d'une ampleur colossale, couvrant un continent entier. Tenter d'analyser cette carte comme une entité monolithique serait une tâche herculéenne. Il est plus judicieux de la subdiviser en régions distinctes : les plaines de l'ouest, les montagnes du centre, les côtes de l'est. Chaque région, bien que faisant partie du tout, peut être étudiée avec ses propres spécificités, ses propres outils d'analyse, avant que les résultats ne soient réintégrés dans une compréhension globale. Cette approche modulaire est la philosophie des \textbf{matrices par blocs}. Elle permet de décomposer un problème de grande dimension, représenté par une matrice massive, en une collection de sous-problèmes plus petits et plus gérables, chacun étant encapsulé dans un "bloc". Les interactions entre ces blocs sont alors étudiées, permettant de reconstruire la solution du problème original avec une efficacité et une clarté accrues.

L'invention de ces concepts n'est pas le fruit du hasard, mais une réponse directe à la nécessité de simplifier l'analyse et le calcul. Certaines transformations géométriques ou physiques, qui apparaissent complexes dans une base arbitraire, se révèlent triviales dans une base spécialement choisie. Les matrices par blocs, quant à elles, sont une manifestation de la stratégie "diviser pour régner", essentielle pour aborder les systèmes de grande taille, qu'ils soient numériques, physiques ou informatiques.

\subsection*{2. Formalisation}

Le passage de l'intuition à la rigueur exige une définition précise des objets et des opérations.

\subsubsection*{A. Définitions Formelles}

Soit $E$ un $\mathbb{K}$-espace vectoriel de dimension finie $n$, où $\mathbb{K}$ désigne un corps commutatif (par exemple, $\mathbb{R}$ ou $\mathbb{C}$).

1.  \textbf{Matrice de passage ($P_{\mathcal{B} \to \mathcal{B}'}$) :}
    Soient $\mathcal{B} = (e_1, e_2, \dots, e_n)$ et $\mathcal{B}' = (e'_1, e'_2, \dots, e'_n)$ deux bases ordonnées de l'espace vectoriel $E$.
    La matrice de passage de la base $\mathcal{B}$ à la base $\mathcal{B}'$, notée $P_{\mathcal{B} \to \mathcal{B}'}$, est la matrice dont les colonnes sont les coordonnées des vecteurs de la nouvelle base $\mathcal{B}'$ exprimées dans l'ancienne base $\mathcal{B}$.
    Formellement, si pour tout $j \in \{1, \dots, n\}$, le vecteur $e'_j$ s'écrit comme une combinaison linéaire des vecteurs de $\mathcal{B}$ :
    $$e'_j = \sum_{i=1}^{n} p_{ij} e_i$$
    alors la matrice $P_{\mathcal{B} \to \mathcal{B}'}$ est donnée par $P_{\mathcal{B} \to \mathcal{B}'} = (p_{ij})_{1 \le i,j \le n}$.
    Cette matrice est également la matrice de l'application identité $Id_E: E \to E$ relativement aux bases $\mathcal{B}'$ au départ et $\mathcal{B}$ à l'arrivée :
    $$P_{\mathcal{B} \to \mathcal{B}'} = \text{Mat}_{\mathcal{B}', \mathcal{B}}(Id_E)$$
    Puisque $\mathcal{B}$ et $\mathcal{B}'$ sont des bases, la matrice $P_{\mathcal{B} \to \mathcal{B}'}$ est toujours inversible.

2.  \textbf{Matrices par blocs :}
    Une matrice $M \in \mathcal{M}_{m,p}(\mathbb{K})$ peut être décomposée en sous-matrices, appelées blocs, en partitionnant ses lignes et ses colonnes.
    Par exemple, une matrice $M$ peut être écrite sous la forme :
    $$M = \begin{pmatrix} A & B \\ C & D \end{pmatrix}$$
    où $A \in \mathcal{M}_{m_1, p_1}(\mathbb{K})$, $B \in \mathcal{M}_{m_1, p_2}(\mathbb{K})$, $C \in \mathcal{M}_{m_2, p_1}(\mathbb{K})$, et $D \in \mathcal{M}_{m_2, p_2}(\mathbb{K})$, avec $m_1+m_2=m$ et $p_1+p_2=p$.
    Les opérations matricielles (somme, produit) peuvent être effectuées sur ces blocs comme s'ils étaient des scalaires, à condition que les dimensions des blocs soient compatibles pour les opérations correspondantes.

\subsubsection*{B. Théorèmes, Propositions \& Lemmes}

\textbf{Théorème 1 (Formule du changement de base pour un vecteur) :}
Soit $u \in E$ un vecteur. Soit $X = \text{Coord}_{\mathcal{B}}(u)$ la colonne des coordonnées de $u$ dans la base $\mathcal{B}$, et $X' = \text{Coord}_{\mathcal{B}'}(u)$ la colonne des coordonnées de $u$ dans la base $\mathcal{B}'$. Alors la relation entre ces coordonnées est donnée par :
$$X = P_{\mathcal{B} \to \mathcal{B}'} X'$$
Ceci signifie que pour obtenir les coordonnées de $u$ dans l'ancienne base $\mathcal{B}$ à partir de ses coordonnées dans la nouvelle base $\mathcal{B}'$, on multiplie par la matrice de passage de $\mathcal{B}$ à $\mathcal{B}'$.

\textbf{Théorème 2 (Formule du changement de base pour un endomorphisme) :}
Soit $f \in \mathcal{L}(E)$ un endomorphisme de $E$. Soit $M = \text{Mat}_{\mathcal{B}}(f)$ la matrice de $f$ dans la base $\mathcal{B}$, et $M' = \text{Mat}_{\mathcal{B}'}(f)$ la matrice de $f$ dans la base $\mathcal{B}'$. Soit $P = P_{\mathcal{B} \to \mathcal{B}'}$ la matrice de passage de $\mathcal{B}$ à $\mathcal{B}'$. Alors la relation entre $M$ et $M'$ est donnée par :
$$M' = P^{-1} M P$$
Les matrices $M$ et $M'$ sont dites \textbf{semblables}. La relation de similitude est une relation d'équivalence.

\subsection*{3. Démonstrations}

Chaque assertion mathématique requiert une preuve rigoureuse, dénuée de toute ambiguïté ou raccourci.

\subsubsection*{Démonstration du Théorème 2 : Formule du changement de base pour un endomorphisme}

Nous souhaitons établir l'égalité $M' = P^{-1} M P$.

1.  \textbf{Initialisation et Cadre des Objets :}
    *   Soit $u$ un vecteur arbitraire de l'espace vectoriel $E$.
    *   Soit $v = f(u)$ l'image de $u$ par l'endomorphisme $f$.
    *   Soient $X \in \mathcal{M}_{n,1}(\mathbb{K})$ la colonne des coordonnées de $u$ dans la base $\mathcal{B}$.
    *   Soient $Y \in \mathcal{M}_{n,1}(\mathbb{K})$ la colonne des coordonnées de $v$ dans la base $\mathcal{B}$.
    *   Soient $X' \in \mathcal{M}_{n,1}(\mathbb{K})$ la colonne des coordonnées de $u$ dans la base $\mathcal{B}'$.
    *   Soient $Y' \in \mathcal{M}_{n,1}(\mathbb{K})$ la colonne des coordonnées de $v$ dans la base $\mathcal{B}'$.
    *   Par définition de la matrice d'un endomorphisme dans une base donnée, nous avons les relations suivantes :
        *   $Y = M X$ (1) : les coordonnées de $v$ dans $\mathcal{B}$ sont obtenues en appliquant $M$ aux coordonnées de $u$ dans $\mathcal{B}$.
        *   $Y' = M' X'$ (2) : les coordonnées de $v$ dans $\mathcal{B}'$ sont obtenues en appliquant $M'$ aux coordonnées de $u$ dans $\mathcal{B}'$.
    *   Par définition de la matrice de passage $P = P_{\mathcal{B} \to \mathcal{B}'}$ (Théorème 1), nous avons les relations suivantes :
        *   $X = P X'$ (3) : les coordonnées de $u$ dans $\mathcal{B}$ sont obtenues à partir de ses coordonnées dans $\mathcal{B}'$ par multiplication par $P$.
        *   $Y = P Y'$ (4) : les coordonnées de $v$ dans $\mathcal{B}$ sont obtenues à partir de ses coordonnées dans $\mathcal{B}'$ par multiplication par $P$.

2.  \textbf{Substitution des Coordonnées de $v$ :}
    Nous partons de l'équation (1) :
    $Y = M X$
    Nous substituons l'expression de $Y$ donnée par l'équation (4) dans cette égalité :
    $P Y' = M X$

3.  \textbf{Substitution des Coordonnées de $u$ :}
    Nous substituons ensuite l'expression de $X$ donnée par l'équation (3) dans l'égalité obtenue à l'étape précédente :
    $P Y' = M (P X')$
    Par associativité de la multiplication matricielle, nous pouvons écrire :
    $P Y' = (M P) X'$

4.  \textbf{Isolement de $Y'$ :}
    Puisque $P$ est une matrice de passage entre deux bases, elle est nécessairement inversible. Nous pouvons donc multiplier l'égalité à gauche par $P^{-1}$ :
    $P^{-1} (P Y') = P^{-1} (M P X')$
    En utilisant l'associativité de la multiplication matricielle et la propriété de l'inverse ($P^{-1} P = I_n$, où $I_n$ est la matrice identité de taille $n \times n$) :
    $(P^{-1} P) Y' = (P^{-1} M P) X'$
    $I_n Y' = (P^{-1} M P) X'$
    $Y' = (P^{-1} M P) X'$

5.  \textbf{Conclusion par Unicité :}
    Nous avons obtenu l'expression $Y' = (P^{-1} M P) X'$.
    Or, par définition de la matrice $M'$ de l'endomorphisme $f$ dans la base $\mathcal{B}'$, nous avons $Y' = M' X'$ (équation (2)).
    Puisque cette égalité $Y' = (P^{-1} M P) X'$ doit être vraie pour tout vecteur $u \in E$ (et donc pour toute colonne de coordonnées $X'$ dans $\mathcal{B}'$), et que la matrice associée à un endomorphisme dans une base donnée est unique, nous en déduisons par identification :
    $M' = P^{-1} M P$.
    La démonstration est achevée.

\subsection*{4. Exercices d'Application}

La maîtrise des concepts s'acquiert par la pratique rigoureuse.

\subsubsection*{Exercice 1 : Application Directe (Coordonnées dans une nouvelle base)}

\textbf{Énoncé :}
Dans l'espace vectoriel $\mathbb{R}^2$ muni de la base canonique $\mathcal{B} = (e_1, e_2)$, où $e_1 = \begin{pmatrix} 1 \\ 0 \end{pmatrix}$ et $e_2 = \begin{pmatrix} 0 \\ 1 \end{pmatrix}$, on considère une nouvelle base $\mathcal{B}' = (e'_1, e'_2)$ définie par les vecteurs $e'_1 = \begin{pmatrix} 1 \\ 1 \end{pmatrix}$ et $e'_2 = \begin{pmatrix} 1 \\ -1 \end{pmatrix}$.
Soit $u$ un vecteur de $\mathbb{R}^2$ dont les coordonnées dans la base canonique $\mathcal{B}$ sont $X = \begin{pmatrix} 4 \\ 2 \end{pmatrix}$.
Calculer les coordonnées $X'$ du vecteur $u$ dans la base $\mathcal{B}'$.

\textbf{Correction Détaillée :}

1.  \textbf{Construction de la matrice de passage $P_{\mathcal{B} \to \mathcal{B}'}$ :}
    Par définition, les colonnes de la matrice de passage $P_{\mathcal{B} \to \mathcal{B}'}$ sont les coordonnées des vecteurs de la nouvelle base $\mathcal{B}'$ exprimées dans l'ancienne base $\mathcal{B}$.
    Les vecteurs $e'_1$ et $e'_2$ sont déjà donnés en coordonnées dans la base canonique $\mathcal{B}$.
    Ainsi, $e'_1 = 1 \cdot e_1 + 1 \cdot e_2$ et $e'_2 = 1 \cdot e_1 + (-1) \cdot e_2$.
    La matrice de passage est donc :
    $$P_{\mathcal{B} \to \mathcal{B}'} = \begin{pmatrix} 1 & 1 \\ 1 & -1 \end{pmatrix}$$

2.  \textbf{Formule de changement de coordonnées :}
    Le Théorème 1 établit la relation $X = P_{\mathcal{B} \to \mathcal{B}'} X'$.
    Pour trouver $X'$, nous devons inverser cette relation : $X' = (P_{\mathcal{B} \to \mathcal{B}'})^{-1} X$.

3.  \textbf{Calcul de l'inverse de la matrice de passage $(P_{\mathcal{B} \to \mathcal{B}'})^{-1}$ :}
    Pour une matrice $A = \begin{pmatrix} a & b \\ c & d \end{pmatrix}$, son déterminant est $\det(A) = ad - bc$. Si $\det(A) \ne 0$, alors $A^{-1} = \frac{1}{\det(A)} \begin{pmatrix} d & -b \\ -c & a \end{pmatrix}$.
    Calculons le déterminant de $P_{\mathcal{B} \to \mathcal{B}'}$ :
    $\det(P_{\mathcal{B} \to \mathcal{B}'}) = (1) \cdot (-1) - (1) \cdot (1) = -1 - 1 = -2$.
    Puisque $\det(P_{\mathcal{B} \to \mathcal{B}'}) = -2 \ne 0$, la matrice est inversible.
    Calculons son inverse :
    $$(P_{\mathcal{B} \to \mathcal{B}'})^{-1} = \frac{1}{-2} \begin{pmatrix} -1 & -1 \\ -1 & 1 \end{pmatrix} = \begin{pmatrix} \frac{-1}{-2} & \frac{-1}{-2} \\ \frac{-1}{-2} & \frac{1}{-2} \end{pmatrix} = \begin{pmatrix} 0.5 & 0.5 \\ 0.5 & -0.5 \end{pmatrix}$$

4.  \textbf{Calcul des coordonnées $X'$ :}
    Nous appliquons la formule $X' = (P_{\mathcal{B} \to \mathcal{B}'})^{-1} X$ :
    $$X' = \begin{pmatrix} 0.5 & 0.5 \\ 0.5 & -0.5 \end{pmatrix} \begin{pmatrix} 4 \\ 2 \end{pmatrix}$$
    Effectuons la multiplication matricielle :
    $$X' = \begin{pmatrix} (0.5 \cdot 4) + (0.5 \cdot 2) \\ (0.5 \cdot 4) + (-0.5 \cdot 2) \end{pmatrix} = \begin{pmatrix} 2 + 1 \\ 2 - 1 \end{pmatrix} = \begin{pmatrix} 3 \\ 1 \end{pmatrix}$$

\textbf{Conclusion :}
Les coordonnées du vecteur $u$ dans la base $\mathcal{B}'$ sont $X' = \begin{pmatrix} 3 \\ 1 \end{pmatrix}$. Cela signifie que $u = 3e'_1 + 1e'_2$.
Vérification : $3e'_1 + 1e'_2 = 3 \begin{pmatrix} 1 \\ 1 \end{pmatrix} + 1 \begin{pmatrix} 1 \\ -1 \end{pmatrix} = \begin{pmatrix} 3 \\ 3 \end{pmatrix} + \begin{pmatrix} 1 \\ -1 \end{pmatrix} = \begin{pmatrix} 3+1 \\ 3-1 \end{pmatrix} = \begin{pmatrix} 4 \\ 2 \end{pmatrix}$, ce qui correspond bien aux coordonnées de $u$ dans la base canonique.

\subsubsection*{Exercice 2 : Niveau Avancé (Matrice par blocs et Inverse)}

\textbf{Énoncé :}
Soit $M$ une matrice carrée de taille $(n+m) \times (n+m)$ sur le corps $\mathbb{K}$, décomposée en blocs de la manière suivante :
$$M = \begin{pmatrix} A & B \\ 0_{m,n} & D \end{pmatrix}$$
où $A \in \mathcal{M}_{n,n}(\mathbb{K})$ est une matrice carrée de taille $n$, $B \in \mathcal{M}_{n,m}(\mathbb{K})$ est une matrice rectangulaire de taille $n \times m$, $0_{m,n} \in \mathcal{M}_{m,n}(\mathbb{K})$ est la matrice nulle de taille $m \times n$, et $D \in \mathcal{M}_{m,m}(\mathbb{K})$ est une matrice carrée de taille $m$.
On suppose que les matrices $A$ et $D$ sont toutes deux inversibles.
Démontrer que la matrice $M$ est inversible et exprimer son inverse $M^{-1}$ sous forme de blocs.

\textbf{Correction Détaillée :}

1.  \textbf{Hypothèse d'inversibilité et forme de l'inverse :}
    Nous cherchons une matrice $M^{-1}$ de même taille que $M$, décomposée en blocs, telle que $M M^{-1} = I_{n+m}$, où $I_{n+m}$ est la matrice identité de taille $(n+m) \times (n+m)$.
    Soit $M^{-1}$ la matrice que nous cherchons, décomposée en blocs de dimensions compatibles :
    $$M^{-1} = \begin{pmatrix} X & Y \\ Z & W \end{pmatrix}$$
    où $X \in \mathcal{M}_{n,n}(\mathbb{K})$, $Y \in \mathcal{M}_{n,m}(\mathbb{K})$, $Z \in \mathcal{M}_{m,n}(\mathbb{K})$, et $W \in \mathcal{M}_{m,m}(\mathbb{K})$.
    La matrice identité $I_{n+m}$ peut également être décomposée en blocs :
    $$I_{n+m} = \begin{pmatrix} I_n & 0_{n,m} \\ 0_{m,n} & I_m \end{pmatrix}$$

2.  \textbf{Calcul du produit $M M^{-1}$ par blocs :}
    Nous effectuons la multiplication matricielle par blocs :
    $$\begin{pmatrix} A & B \\ 0_{m,n} & D \end{pmatrix} \begin{pmatrix} X & Y \\ Z & W \end{pmatrix} = \begin{pmatrix} (A X + B Z) & (A Y + B W) \\ (0_{m,n} X + D Z) & (0_{m,n} Y + D W) \end{pmatrix}$$
    Simplifions les termes impliquant la matrice nulle :
    $$M M^{-1} = \begin{pmatrix} A X + B Z & A Y + B W \\ D Z & D W \end{pmatrix}$$

3.  \textbf{Identification avec la matrice identité par blocs :}
    Nous égalons le résultat du produit $M M^{-1}$ à la matrice identité $I_{n+m}$ :
    $$\begin{pmatrix} A X + B Z & A Y + B W \\ D Z & D W \end{pmatrix} = \begin{pmatrix} I_n & 0_{n,m} \\ 0_{m,n} & I_m \end{pmatrix}$$
    Cette égalité matricielle par blocs conduit à un système de quatre équations matricielles :
    (a) $A X + B Z = I_n$
    (b) $A Y + B W = 0_{n,m}$
    (c) $D Z = 0_{m,n}$
    (d) $D W = I_m$

4.  \textbf{Résolution du système d'équations pour trouver $X, Y, Z, W$ :}

    *   \textbf{À partir de l'équation (c) :} $D Z = 0_{m,n}$
        Puisque $D$ est une matrice inversible (par hypothèse), nous pouvons multiplier cette équation à gauche par $D^{-1}$ :
        $D^{-1} (D Z) = D^{-1} 0_{m,n}$
        $(D^{-1} D) Z = 0_{m,n}$
        $I_m Z = 0_{m,n}$
        $Z = 0_{m,n}$

    *   \textbf{À partir de l'équation (d) :} $D W = I_m$
        Puisque $D$ est inversible, nous pouvons multiplier cette équation à gauche par $D^{-1}$ :
        $D^{-1} (D W) = D^{-1} I_m$
        $(D^{-1} D) W = D^{-1}$
        $I_m W = D^{-1}$
        $W = D^{-1}$

    *   \textbf{À partir de l'équation (a) :} $A X + B Z = I_n$
        Nous substituons la valeur de $Z = 0_{m,n}$ que nous venons de trouver :
        $A X + B (0_{m,n}) = I_n$
        $A X + 0_{n,n} = I_n$
        $A X = I_n$
        Puisque $A$ est une matrice inversible (par hypothèse), nous pouvons multiplier cette équation à gauche par $A^{-1}$ :
        $A^{-1} (A X) = A^{-1} I_n$
        $(A^{-1} A) X = A^{-1}$
        $I_n X = A^{-1}$
        $X = A^{-1}$

    *   \textbf{À partir de l'équation (b) :} $A Y + B W = 0_{n,m}$
        Nous substituons la valeur de $W = D^{-1}$ que nous avons trouvée :
        $A Y + B D^{-1} = 0_{n,m}$
        Nous isolons le terme $A Y$ :
        $A Y = -B D^{-1}$
        Puisque $A$ est une matrice inversible, nous pouvons multiplier cette équation à gauche par $A^{-1}$ :
        $A^{-1} (A Y) = A^{-1} (-B D^{-1})$
        $(A^{-1} A) Y = -A^{-1} B D^{-1}$
        $I_n Y = -A^{-1} B D^{-1}$
        $Y = -A^{-1} B D^{-1}$

5.  \textbf{Conclusion :}
    Nous avons trouvé des expressions uniques pour tous les blocs $X, Y, Z, W$. Cela démontre que la matrice $M$ est inversible, et son inverse $M^{-1}$ est donnée par :
    $$M^{-1} = \begin{pmatrix} A^{-1} & -A^{-1} B D^{-1} \\ 0_{m,n} & D^{-1} \end{pmatrix}$$

\subsection*{5. Application en Intelligence Artificielle}

Les concepts de changement de base et de matrices par blocs ne sont pas de simples abstractions mathématiques ; ils constituent le socle de nombreuses techniques fondamentales en intelligence artificielle, permettant de transformer des données brutes en représentations plus significatives et plus efficaces pour l'apprentissage et l'analyse.

Le \textbf{changement de base} est au cœur de l'\textbf{extraction de caractéristiques (Feature Engineering)} et de la \textbf{réduction de dimensionnalité}. Dans de nombreux problèmes d'apprentissage automatique, les données initiales (par exemple, les valeurs de pixels d'une image, les mots d'un texte) résident dans un espace de très haute dimension, où les motifs pertinents sont souvent noyés dans le bruit ou exprimés de manière redondante. En changeant de base, on projette ces données dans un nouvel espace où les caractéristiques discriminantes sont amplifiées, les corrélations sont simplifiées, et la dimensionnalité peut être réduite sans perte significative d'information. C'est une transformation du "point de vue" sur les données pour en révéler la structure intrinsèque.

Un exemple emblématique de cette application est la \textbf{compression d'images JPEG}, qui repose sur la \textbf{Transformée en Cosinus Discrète (DCT)}. La DCT est un changement de base orthogonal qui transforme un signal (ici, les valeurs de pixels) du domaine spatial (où chaque pixel a une position) vers le domaine fréquentiel (où chaque coefficient représente l'amplitude d'une certaine fréquence spatiale).
Le processus se déroule comme suit :
1.  \textbf{Découpage par blocs :} L'image est d'abord divisée en petits blocs de pixels, généralement de taille $8 \times 8$. Cette étape illustre parfaitement l'utilisation des \textbf{matrices par blocs} : chaque bloc est traité indépendamment, ce qui permet de gérer la complexité d'une image entière en la décomposant en sous-problèmes.
2.  \textbf{Application de la DCT :} Pour chaque bloc $8 \times 8$, une transformation de base est appliquée. La base de la DCT est composée de fonctions cosinus de différentes fréquences et orientations. Dans cette nouvelle base, l'information visuelle est "compactée" : les coefficients correspondant aux basses fréquences (qui représentent les grandes structures et les couleurs générales) ont généralement des amplitudes élevées, tandis que les coefficients des hautes fréquences (qui représentent les détails fins et les textures) ont des amplitudes plus faibles.
3.  \textbf{Quantification et compression :} C'est là que la compression a lieu. Les coefficients de haute fréquence, qui contribuent moins à la perception visuelle humaine et sont souvent associés au bruit, sont quantifiés de manière plus grossière (c'est-à-dire que leur précision est réduite, voire qu'ils sont mis à zéro). Cette opération est possible et efficace précisément parce que le changement de base a regroupé l'information essentielle dans un petit nombre de coefficients.
4.  \textbf{Encodage :} Les coefficients quantifiés sont ensuite encodés de manière efficace.

Ainsi, la compression JPEG est une démonstration éloquente de la synergie entre le changement de base (DCT) et les matrices par blocs (traitement par $8 \times 8$ blocs). Elle permet de réduire drastiquement la taille des fichiers images en éliminant les informations redondantes ou moins perceptibles, tout en préservant une qualité visuelle acceptable, en exploitant la capacité des changements de base à révéler la structure essentielle des données.

\subsection*{6. Liens Sémantiques}

Les concepts abordés ici s'inscrivent dans une progression logique de l'algèbre linéaire et sont des prérequis indispensables à des notions plus avancées.

*   \textbf{Concepts Précédents requis :}
    *   [[Jalon-7.md|Jalon 7 (Espaces vectoriels abstraits)]] : Compréhension des espaces vectoriels, des bases et des coordonnées.
    *   [[Jalon-9.md|Jalon 9 (Applications linéaires et matrices)]] : Maîtrise des applications linéaires, de leur représentation matricielle et des opérations matricielles.

*   \textbf{Concepts Futurs dépendants :}
    *   [[Jalon 29 (Éléments propres)]] : La diagonalisation des matrices, qui est un cas particulier de changement de base où la matrice de l'endomorphisme devient diagonale, est fondamentale pour l'étude des valeurs et vecteurs propres.
    *   [[Jalon 30 (Trigonalisation d'endomorphismes et décomposition de Dunford.)]] : Ces techniques visent à simplifier la matrice d'un endomorphisme par un changement de base, même lorsque la diagonalisation n'est pas possible.
    *   [[Jalon 80 (Transformée de Fourier dans L^1)]] : La transformée de Fourier est une généralisation du concept de changement de base à des espaces de fonctions, transformant un signal du domaine temporel au domaine fréquentiel, essentielle en traitement du signal et en physique.

\subsection*{Genèse Narrative}

Le concept de changement de base...

\subsection*{Énoncé Symbolique Strict}

Soit $E$ un espace vectoriel...
\section*{Illustration Visuelle (TikZ)}
\begin{center}
\begin{tikzpicture}
\draw[->] (0,0) -- (2,0) node[right] {$e_1$};
\draw[->] (0,0) -- (0,2) node[above] {$e_2$};
\draw[->, thick, red] (0,0) -- (1,2) node[above right] {$e'_1$};
\draw[->, thick, red] (0,0) -- (-1,1) node[above left] {$e'_2$};
\end{tikzpicture}
\end{center}
\chapter{Exercices d'Application}
\section*{Exercice 01}
\subsection*{Énoncé}
Soit $E = \mathbb{R}^2$ muni de sa base canonique $\mathcal{B} = (e_1, e_2)$.
On considère les vecteurs $e'_1 = (1, 2)$ et $e'_2 = (-1, 1)$.
1. Montrer que $\mathcal{B}' = (e'_1, e'_2)$ forme une base de $E$.
2. Écrire la matrice de passage $P$ de la base $\mathcal{B}$ à la base $\mathcal{B}'$.
3. Calculer la matrice $P^{-1}$.
4. Soit un vecteur $u$ de $E$ dont les coordonnées dans la base $\mathcal{B}$ sont $X = \begin{pmatrix} 2 \\ 5 \end{pmatrix}$.
   Calculer ses coordonnées $X'$ dans la base $\mathcal{B}'$ en utilisant $P^{-1}$.

\subsection*{Correction}
\textbf{1. Montrer que $\mathcal{B}'$ est une base de $E$ :}
Puisque $\dim(E) = 2$ et que la famille $\mathcal{B}'$ contient $2$ vecteurs, il suffit de vérifier qu'elle est libre.
Soient $\lambda_1, \lambda_2 \in \mathbb{R}$ tels que $\lambda_1 e'_1 + \lambda_2 e'_2 = 0_E$.
$\lambda_1(1, 2) + \lambda_2(-1, 1) = (0, 0)$
On obtient le système :
$\begin{cases} \lambda_1 - \lambda_2 = 0 \\ 2\lambda_1 + \lambda_2 = 0 \end{cases}$
De la première équation, on a $\lambda_1 = \lambda_2$.
En remplaçant dans la seconde : $2\lambda_1 + \lambda_1 = 3\lambda_1 = 0 \implies \lambda_1 = 0$.
D'où $\lambda_2 = 0$. La famille est libre, c'est donc une base de $\mathbb{R}^2$.

\textbf{2. Matrice de passage $P$ :}
Les colonnes de $P$ sont les coordonnées des vecteurs de la nouvelle base $\mathcal{B}'$ dans l'ancienne base $\mathcal{B}$.
$e'_1 = 1e_1 + 2e_2 \implies \text{colonne 1 } = \begin{pmatrix} 1 \\ 2 \end{pmatrix}$
$e'_2 = -1e_1 + 1e_2 \implies \text{colonne 2 } = \begin{pmatrix} -1 \\ 1 \end{pmatrix}$
$P = \begin{pmatrix} 1 & -1 \\ 2 & 1 \end{pmatrix}$.

\textbf{3. Calcul de $P^{-1}$ :}
Le déterminant de $P$ est $\det(P) = (1)(1) - (-1)(2) = 1 + 2 = 3$.
La formule de l'inverse d'une matrice $2 \times 2$ donne :
$P^{-1} = \frac{1}{\det(P)} \begin{pmatrix} d & -b \\ -c & a \end{pmatrix} = \frac{1}{3} \begin{pmatrix} 1 & 1 \\ -2 & 1 \end{pmatrix} = \begin{pmatrix} 1/3 & 1/3 \\ -2/3 & 1/3 \end{pmatrix}$.

\textbf{4. Coordonnées de $u$ dans $\mathcal{B}'$ :}
On utilise la formule $X' = P^{-1} X$.
$X' = \frac{1}{3} \begin{pmatrix} 1 & 1 \\ -2 & 1 \end{pmatrix} \begin{pmatrix} 2 \\ 5 \end{pmatrix} = \frac{1}{3} \begin{pmatrix} 1 \times 2 + 1 \times 5 \\ -2 \times 2 + 1 \times 5 \end{pmatrix} = \frac{1}{3} \begin{pmatrix} 7 \\ 1 \end{pmatrix} = \begin{pmatrix} 7/3 \\ 1/3 \end{pmatrix}$.
Les coordonnées de $u$ dans la base $\mathcal{B}'$ sont donc $(7/3, 1/3)$.







\subsection*{Correction exhaustive (Protocole d'Exégèse)}

\textbf{1. Énoncé symbolique et Typage Chirurgical :}
Soit $E$ un $\mathbb{K}$-espace vectoriel de dimension finie $n$. Soient $\mathcal{B} = (e_1, \dots, e_n)$ et $\mathcal{B}' = (e'_1, \dots, e'_n)$ deux bases de $E$.
La matrice de passage $P = \operatorname{Mat}_{\mathcal{B}}(\mathcal{B}')$ a pour colonnes les composantes des vecteurs $e'_j$ exprimées dans la base $\mathcal{B}$.
Soit $u \in E$, on note $X = \operatorname{Mat}_{\mathcal{B}}(u)$ et $X' = \operatorname{Mat}_{\mathcal{B}'}(u)$.

\textbf{2. Démonstration de l'inversibilité de la matrice de passage (Zéro ellipse) :}
Puisque $\mathcal{B}'$ est une base, l'application linéaire $f : E \to E$ définie par $f(e_j) = e'_j$ pour tout $j \in \{1, \dots, n\}$ transforme une base en une base, c'est donc un automorphisme de $E$.
Or, par définition, $P = \operatorname{Mat}_{\mathcal{B}}(f)$.
Puisque $f$ est bijective, son déterminant est non nul : $\det(P) \neq 0$.
Ainsi, $P \in \operatorname{GL}_n(\mathbb{K})$ est inversible.

\textbf{3. Calcul de la relation de changement de base :}
On sait que $u = \sum_{j=1}^n x'_j e'_j$.
Remplaçons $e'_j$ par son expression dans $\mathcal{B}$ : $e'_j = \sum_{i=1}^n P_{ij} e_i$.
Alors $u = \sum_{j=1}^n x'_j \left( \sum_{i=1}^n P_{ij} e_i \right) = \sum_{i=1}^n \left( \sum_{j=1}^n P_{ij} x'_j \right) e_i$.
Puisque la décomposition de $u$ dans $\mathcal{B}$ est unique, on a $x_i = \sum_{j=1}^n P_{ij} x'_j$.
En notation matricielle, ceci s'écrit exactement : $X = P X'$.
En multipliant à gauche par $P^{-1}$, on obtient $X' = P^{-1} X$.

\textbf{4. Exemples et Cas Pathologiques :}
- *Cas trivial :* $\mathcal{B} = \mathcal{B}' \implies P = I_n$ (matrice identité).
- *Contre-exemple (famille liée) :* Si $\mathcal{B}'$ n'était pas libre, $P$ ne serait pas inversible et la relation $X = P X'$ ne pourrait pas être inversée, signifiant que les coordonnées dans $\mathcal{B}'$ ne seraient pas uniquement déterminées.


\section*{Exercice 02}
\subsection*{Énoncé}
Soit $E = \mathbb{R}^3$ muni de sa base canonique $\mathcal{B} = (e_1, e_2, e_3)$.
On considère les vecteurs $v_1 = (1, 1, 0)$, $v_2 = (0, 1, 1)$ et $v_3 = (1, 0, 1)$.
1. Prouver que $\mathcal{B}' = (v_1, v_2, v_3)$ est une base de $E$.
2. Donner la matrice de passage $P$ de $\mathcal{B}$ à $\mathcal{B}'$.
3. Calculer l'inverse de la matrice $P$ par la méthode du pivot de Gauss ou des cofacteurs.
4. On considère le vecteur $w$ dont les coordonnées dans la base $\mathcal{B}'$ sont $X' = \begin{pmatrix} 1 \\ -2 \\ 3 \end{pmatrix}$. Quelles sont ses coordonnées $X$ dans la base $\mathcal{B}$ ?

\subsection*{Correction}
\textbf{1. Prouver que $\mathcal{B}'$ est une base :}
La famille $\mathcal{B}'$ comporte 3 vecteurs dans $\mathbb{R}^3$ de dimension 3. Il suffit de vérifier qu'elle est libre.
Soit la combinaison linéaire $\lambda_1 v_1 + \lambda_2 v_2 + \lambda_3 v_3 = (0,0,0)$.
Cela se traduit par le système :
$\begin{cases} \lambda_1 + \lambda_3 = 0 \quad (L_1)\\ \lambda_1 + \lambda_2 = 0 \quad (L_2)\\ \lambda_2 + \lambda_3 = 0 \quad (L_3) \end{cases}$
De $(L_1)$ on a $\lambda_3 = -\lambda_1$.
De $(L_2)$ on a $\lambda_2 = -\lambda_1$.
En remplaçant dans $(L_3)$ : $(-\lambda_1) + (-\lambda_1) = 0 \implies -2\lambda_1 = 0 \implies \lambda_1 = 0$.
On en déduit que $\lambda_2 = \lambda_3 = 0$. La famille est libre, c'est une base.

\textbf{2. Matrice de passage $P$ :}
Les colonnes de $P$ sont formées par les vecteurs $v_1, v_2, v_3$ exprimés dans la base canonique.
$P = \begin{pmatrix} 1 & 0 & 1 \\ 1 & 1 & 0 \\ 0 & 1 & 1 \end{pmatrix}$.

\textbf{3. Calcul de $P^{-1}$ :}
Utilisons la méthode des cofacteurs.
$\det(P) = 1(1 \cdot 1 - 0 \cdot 1) - 0 + 1(1 \cdot 1 - 1 \cdot 0) = 1 + 1 = 2$.
La matrice des cofacteurs (comatrice) est :
$C = \begin{pmatrix} +(1-0) & -(1-0) & +(1-0) \\ -(0-1) & +(1-0) & -(1-0) \\ +(0-1) & -(0-1) & +(1-0) \end{pmatrix} = \begin{pmatrix} 1 & -1 & 1 \\ 1 & 1 & -1 \\ -1 & 1 & 1 \end{pmatrix}$.
L'inverse est donné par $P^{-1} = \frac{1}{\det P} C^T$.
$P^{-1} = \frac{1}{2} \begin{pmatrix} 1 & 1 & -1 \\ -1 & 1 & 1 \\ 1 & -1 & 1 \end{pmatrix}$.

\textbf{4. Coordonnées de $w$ dans $\mathcal{B}$ :}
La relation entre les anciennes et nouvelles coordonnées est donnée par la formule $X = P X'$.
$X = \begin{pmatrix} 1 & 0 & 1 \\ 1 & 1 & 0 \\ 0 & 1 & 1 \end{pmatrix} \begin{pmatrix} 1 \\ -2 \\ 3 \end{pmatrix} = \begin{pmatrix} 1 \times 1 + 0 \times (-2) + 1 \times 3 \\ 1 \times 1 + 1 \times (-2) + 0 \times 3 \\ 0 \times 1 + 1 \times (-2) + 1 \times 3 \end{pmatrix} = \begin{pmatrix} 1 + 3 \\ 1 - 2 \\ -2 + 3 \end{pmatrix} = \begin{pmatrix} 4 \\ -1 \\ 1 \end{pmatrix}$.
Les coordonnées de $w$ dans la base canonique sont donc $(4, -1, 1)$.







\subsection*{Correction exhaustive (Protocole d'Exégèse)}

\textbf{1. Énoncé symbolique et Typage Chirurgical :}
Soit $E$ un $\mathbb{K}$-espace vectoriel de dimension finie $n$. Soient $\mathcal{B} = (e_1, \dots, e_n)$ et $\mathcal{B}' = (e'_1, \dots, e'_n)$ deux bases de $E$.
La matrice de passage $P = \operatorname{Mat}_{\mathcal{B}}(\mathcal{B}')$ a pour colonnes les composantes des vecteurs $e'_j$ exprimées dans la base $\mathcal{B}$.
Soit $u \in E$, on note $X = \operatorname{Mat}_{\mathcal{B}}(u)$ et $X' = \operatorname{Mat}_{\mathcal{B}'}(u)$.

\textbf{2. Démonstration de l'inversibilité de la matrice de passage (Zéro ellipse) :}
Puisque $\mathcal{B}'$ est une base, l'application linéaire $f : E \to E$ définie par $f(e_j) = e'_j$ pour tout $j \in \{1, \dots, n\}$ transforme une base en une base, c'est donc un automorphisme de $E$.
Or, par définition, $P = \operatorname{Mat}_{\mathcal{B}}(f)$.
Puisque $f$ est bijective, son déterminant est non nul : $\det(P) \neq 0$.
Ainsi, $P \in \operatorname{GL}_n(\mathbb{K})$ est inversible.

\textbf{3. Calcul de la relation de changement de base :}
On sait que $u = \sum_{j=1}^n x'_j e'_j$.
Remplaçons $e'_j$ par son expression dans $\mathcal{B}$ : $e'_j = \sum_{i=1}^n P_{ij} e_i$.
Alors $u = \sum_{j=1}^n x'_j \left( \sum_{i=1}^n P_{ij} e_i \right) = \sum_{i=1}^n \left( \sum_{j=1}^n P_{ij} x'_j \right) e_i$.
Puisque la décomposition de $u$ dans $\mathcal{B}$ est unique, on a $x_i = \sum_{j=1}^n P_{ij} x'_j$.
En notation matricielle, ceci s'écrit exactement : $X = P X'$.
En multipliant à gauche par $P^{-1}$, on obtient $X' = P^{-1} X$.

\textbf{4. Exemples et Cas Pathologiques :}
- *Cas trivial :* $\mathcal{B} = \mathcal{B}' \implies P = I_n$ (matrice identité).
- *Contre-exemple (famille liée) :* Si $\mathcal{B}'$ n'était pas libre, $P$ ne serait pas inversible et la relation $X = P X'$ ne pourrait pas être inversée, signifiant que les coordonnées dans $\mathcal{B}'$ ne seraient pas uniquement déterminées.


\section*{Exercice 03}
\subsection*{Énoncé}
Soit $E$ un $\mathbb{R}$-espace vectoriel de dimension 2. Soit $\mathcal{B} = (e_1, e_2)$ une base de $E$.
Soit $\mathcal{B}' = (e'_1, e'_2)$ la famille définie par :
$e'_1 = 2e_1 + e_2$
$e'_2 = 3e_1 + 2e_2$

1. Montrer que $\mathcal{B}'$ est une base de $E$.
2. Écrire la matrice de passage $P$ de $\mathcal{B}$ à $\mathcal{B}'$.
3. Calculer l'inverse de $P$.
4. Soit $u \in E$ le vecteur de coordonnées $X = \begin{pmatrix} 1 \\ -1 \end{pmatrix}$ dans $\mathcal{B}$. Calculer ses coordonnées $X'$ dans $\mathcal{B}'$.

\subsection*{Correction}
\textbf{1. Montrons que $\mathcal{B}'$ est une base de $E$ :}
Il suffit de montrer que la famille $(e'_1, e'_2)$ est libre, car elle contient 2 vecteurs et $\dim(E) = 2$.
Soient $\lambda_1, \lambda_2 \in \mathbb{R}$ tels que $\lambda_1 e'_1 + \lambda_2 e'_2 = 0_E$.
$\lambda_1 (2e_1 + e_2) + \lambda_2 (3e_1 + 2e_2) = 0_E$
$(2\lambda_1 + 3\lambda_2)e_1 + (\lambda_1 + 2\lambda_2)e_2 = 0_E$
Comme $\mathcal{B} = (e_1, e_2)$ est une base, c'est une famille libre. On obtient le système :
$\begin{cases} 2\lambda_1 + 3\lambda_2 = 0 \\ \lambda_1 + 2\lambda_2 = 0 \end{cases}$
De la deuxième équation, on tire $\lambda_1 = -2\lambda_2$.
En remplaçant dans la première : $2(-2\lambda_2) + 3\lambda_2 = 0 \implies -\lambda_2 = 0 \implies \lambda_2 = 0$.
Puis $\lambda_1 = -2(0) = 0$.
La famille $(e'_1, e'_2)$ est donc libre. C'est bien une base de $E$.

\textbf{2. Matrice de passage $P$ :}
Par définition, la matrice de passage $P_{\mathcal{B} \to \mathcal{B}'}$ contient en colonnes les coordonnées des vecteurs de la nouvelle base $\mathcal{B}'$ exprimés dans l'ancienne base $\mathcal{B}$.
$e'_1 = 2e_1 + 1e_2 \implies \text{colonne 1 } = \begin{pmatrix} 2 \\ 1 \end{pmatrix}$
$e'_2 = 3e_1 + 2e_2 \implies \text{colonne 2 } = \begin{pmatrix} 3 \\ 2 \end{pmatrix}$
Donc $P = \begin{pmatrix} 2 & 3 \\ 1 & 2 \end{pmatrix}$.

\textbf{3. Calcul de $P^{-1}$ :}
Le déterminant de $P$ est $\det(P) = 2 \times 2 - 3 \times 1 = 4 - 3 = 1$.
Puisque $\det(P) \neq 0$, $P$ est bien inversible.
La formule de l'inverse pour une matrice $2 \times 2$ $\begin{pmatrix} a & b \\ c & d \end{pmatrix}$ est $\frac{1}{ad-bc} \begin{pmatrix} d & -b \\ -c & a \end{pmatrix}$.
Ici, $P^{-1} = \frac{1}{1} \begin{pmatrix} 2 & -3 \\ -1 & 2 \end{pmatrix} = \begin{pmatrix} 2 & -3 \\ -1 & 2 \end{pmatrix}$.

\textbf{4. Coordonnées de $u$ dans $\mathcal{B}'$ :}
La formule de changement de coordonnées est $X = P X'$, ce qui équivaut à $X' = P^{-1} X$.
On a $X = \begin{pmatrix} 1 \\ -1 \end{pmatrix}$.
Calculons $X'$ :
$X' = \begin{pmatrix} 2 & -3 \\ -1 & 2 \end{pmatrix} \begin{pmatrix} 1 \\ -1 \end{pmatrix} = \begin{pmatrix} 2 \times 1 + (-3) \times (-1) \\ (-1) \times 1 + 2 \times (-1) \end{pmatrix} = \begin{pmatrix} 2 + 3 \\ -1 - 2 \end{pmatrix} = \begin{pmatrix} 5 \\ -3 \end{pmatrix}$.
Les coordonnées de $u$ dans la base $\mathcal{B}'$ sont donc $\begin{pmatrix} 5 \\ -3 \end{pmatrix}$.







\subsection*{Correction exhaustive (Protocole d'Exégèse)}

\textbf{1. Énoncé symbolique et Typage Chirurgical :}
Soit $E$ un $\mathbb{K}$-espace vectoriel de dimension finie $n$. Soient $\mathcal{B} = (e_1, \dots, e_n)$ et $\mathcal{B}' = (e'_1, \dots, e'_n)$ deux bases de $E$.
La matrice de passage $P = \operatorname{Mat}_{\mathcal{B}}(\mathcal{B}')$ a pour colonnes les composantes des vecteurs $e'_j$ exprimées dans la base $\mathcal{B}$.
Soit $u \in E$, on note $X = \operatorname{Mat}_{\mathcal{B}}(u)$ et $X' = \operatorname{Mat}_{\mathcal{B}'}(u)$.

\textbf{2. Démonstration de l'inversibilité de la matrice de passage (Zéro ellipse) :}
Puisque $\mathcal{B}'$ est une base, l'application linéaire $f : E \to E$ définie par $f(e_j) = e'_j$ pour tout $j \in \{1, \dots, n\}$ transforme une base en une base, c'est donc un automorphisme de $E$.
Or, par définition, $P = \operatorname{Mat}_{\mathcal{B}}(f)$.
Puisque $f$ est bijective, son déterminant est non nul : $\det(P) \neq 0$.
Ainsi, $P \in \operatorname{GL}_n(\mathbb{K})$ est inversible.

\textbf{3. Calcul de la relation de changement de base :}
On sait que $u = \sum_{j=1}^n x'_j e'_j$.
Remplaçons $e'_j$ par son expression dans $\mathcal{B}$ : $e'_j = \sum_{i=1}^n P_{ij} e_i$.
Alors $u = \sum_{j=1}^n x'_j \left( \sum_{i=1}^n P_{ij} e_i \right) = \sum_{i=1}^n \left( \sum_{j=1}^n P_{ij} x'_j \right) e_i$.
Puisque la décomposition de $u$ dans $\mathcal{B}$ est unique, on a $x_i = \sum_{j=1}^n P_{ij} x'_j$.
En notation matricielle, ceci s'écrit exactement : $X = P X'$.
En multipliant à gauche par $P^{-1}$, on obtient $X' = P^{-1} X$.

\textbf{4. Exemples et Cas Pathologiques :}
- *Cas trivial :* $\mathcal{B} = \mathcal{B}' \implies P = I_n$ (matrice identité).
- *Contre-exemple (famille liée) :* Si $\mathcal{B}'$ n'était pas libre, $P$ ne serait pas inversible et la relation $X = P X'$ ne pourrait pas être inversée, signifiant que les coordonnées dans $\mathcal{B}'$ ne seraient pas uniquement déterminées.


\section*{Exercice 04}
\subsection*{Énoncé}
Dans l'espace vectoriel $E = \mathbb{R}^2$, on considère la base canonique $\mathcal{B} = (e_1, e_2)$.
On définit l'endomorphisme $f$ par sa matrice dans $\mathcal{B}$ :
$A = \begin{pmatrix} 1 & 4 \\ 2 & 3 \end{pmatrix}$.
On pose $u_1 = (1, 1)$ et $u_2 = (2, -1)$.
1. Montrer que $\mathcal{B}' = (u_1, u_2)$ est une base de $E$.
2. Écrire la matrice de passage $P$ de $\mathcal{B}$ à $\mathcal{B}'$, et calculer $P^{-1}$.
3. Calculer $f(u_1)$ et $f(u_2)$ sous forme de vecteurs de la base canonique. En déduire que $u_1$ et $u_2$ sont des vecteurs propres de $f$.
4. Quelle est la matrice $A'$ de $f$ dans la base $\mathcal{B}'$ ?
5. Retrouver le résultat précédent en effectuant le produit matriciel $P^{-1} A P$.

\subsection*{Correction}
\textbf{1. Montrer que $\mathcal{B}'$ est une base :}
Il suffit de calculer le déterminant des vecteurs de la famille dans la base canonique :
$\det(u_1, u_2) = \begin{vmatrix} 1 & 2 \\ 1 & -1 \end{vmatrix} = 1(-1) - 1(2) = -1 - 2 = -3$.
Comme le déterminant est non nul, la famille est libre. De cardinal 2 en dimension 2, c'est une base.

\textbf{2. Matrice de passage $P$ et son inverse :}
$P = \begin{pmatrix} 1 & 2 \\ 1 & -1 \end{pmatrix}$.
$P^{-1} = \frac{1}{-3} \begin{pmatrix} -1 & -2 \\ -1 & 1 \end{pmatrix} = \begin{pmatrix} 1/3 & 2/3 \\ 1/3 & -1/3 \end{pmatrix}$.

\textbf{3. Calcul de $f(u_1)$ et $f(u_2)$ :}
Le vecteur coordonnée de $u_1$ est $X_1 = \begin{pmatrix} 1 \\ 1 \end{pmatrix}$.
$f(u_1)$ a pour vecteur coordonnée $AX_1 = \begin{pmatrix} 1 & 4 \\ 2 & 3 \end{pmatrix} \begin{pmatrix} 1 \\ 1 \end{pmatrix} = \begin{pmatrix} 1+4 \\ 2+3 \end{pmatrix} = \begin{pmatrix} 5 \\ 5 \end{pmatrix} = 5 \begin{pmatrix} 1 \\ 1 \end{pmatrix}$.
Donc $f(u_1) = 5u_1$. $u_1$ est un vecteur propre associé à la valeur propre $5$.

Le vecteur coordonnée de $u_2$ est $X_2 = \begin{pmatrix} 2 \\ -1 \end{pmatrix}$.
$f(u_2)$ a pour vecteur coordonnée $AX_2 = \begin{pmatrix} 1 & 4 \\ 2 & 3 \end{pmatrix} \begin{pmatrix} 2 \\ -1 \end{pmatrix} = \begin{pmatrix} 2-4 \\ 4-3 \end{pmatrix} = \begin{pmatrix} -2 \\ 1 \end{pmatrix} = -1 \begin{pmatrix} 2 \\ -1 \end{pmatrix}$.
Donc $f(u_2) = -u_2$. $u_2$ est un vecteur propre associé à la valeur propre $-1$.

\textbf{4. Matrice $A'$ de $f$ dans $\mathcal{B}'$ :}
Les colonnes de $A'$ sont formées par les coordonnées de $f(u_1)$ et $f(u_2)$ dans la base $(u_1, u_2)$.
Comme $f(u_1) = 5u_1 + 0u_2$ et $f(u_2) = 0u_1 - 1u_2$, on a directement :
$A' = \begin{pmatrix} 5 & 0 \\ 0 & -1 \end{pmatrix}$.

\textbf{5. Vérification avec $P^{-1} A P$ :}
Calculons d'abord $AP$ :
$AP = \begin{pmatrix} 1 & 4 \\ 2 & 3 \end{pmatrix} \begin{pmatrix} 1 & 2 \\ 1 & -1 \end{pmatrix} = \begin{pmatrix} 1 \times 1 + 4 \times 1 & 1 \times 2 + 4 \times (-1) \\ 2 \times 1 + 3 \times 1 & 2 \times 2 + 3 \times (-1) \end{pmatrix} = \begin{pmatrix} 5 & -2 \\ 5 & 1 \end{pmatrix}$.
Calculons ensuite $P^{-1}(AP)$ :
$P^{-1}(AP) = \begin{pmatrix} 1/3 & 2/3 \\ 1/3 & -1/3 \end{pmatrix} \begin{pmatrix} 5 & -2 \\ 5 & 1 \end{pmatrix} = \begin{pmatrix} \frac{5}{3} + \frac{10}{3} & \frac{-2}{3} + \frac{2}{3} \\ \frac{5}{3} - \frac{5}{3} & \frac{-2}{3} - \frac{1}{3} \end{pmatrix} = \begin{pmatrix} \frac{15}{3} & 0 \\ 0 & \frac{-3}{3} \end{pmatrix} = \begin{pmatrix} 5 & 0 \\ 0 & -1 \end{pmatrix}$.
On retrouve bien la matrice diagonale $A'$.







\subsection*{Correction exhaustive (Protocole d'Exégèse)}

\textbf{1. Énoncé symbolique et Typage Chirurgical :}
Soit $E$ un $\mathbb{K}$-espace vectoriel de dimension finie $n$. Soient $\mathcal{B} = (e_1, \dots, e_n)$ et $\mathcal{B}' = (e'_1, \dots, e'_n)$ deux bases de $E$.
La matrice de passage $P = \operatorname{Mat}_{\mathcal{B}}(\mathcal{B}')$ a pour colonnes les composantes des vecteurs $e'_j$ exprimées dans la base $\mathcal{B}$.
Soit $u \in E$, on note $X = \operatorname{Mat}_{\mathcal{B}}(u)$ et $X' = \operatorname{Mat}_{\mathcal{B}'}(u)$.

\textbf{2. Démonstration de l'inversibilité de la matrice de passage (Zéro ellipse) :}
Puisque $\mathcal{B}'$ est une base, l'application linéaire $f : E \to E$ définie par $f(e_j) = e'_j$ pour tout $j \in \{1, \dots, n\}$ transforme une base en une base, c'est donc un automorphisme de $E$.
Or, par définition, $P = \operatorname{Mat}_{\mathcal{B}}(f)$.
Puisque $f$ est bijective, son déterminant est non nul : $\det(P) \neq 0$.
Ainsi, $P \in \operatorname{GL}_n(\mathbb{K})$ est inversible.

\textbf{3. Calcul de la relation de changement de base :}
On sait que $u = \sum_{j=1}^n x'_j e'_j$.
Remplaçons $e'_j$ par son expression dans $\mathcal{B}$ : $e'_j = \sum_{i=1}^n P_{ij} e_i$.
Alors $u = \sum_{j=1}^n x'_j \left( \sum_{i=1}^n P_{ij} e_i \right) = \sum_{i=1}^n \left( \sum_{j=1}^n P_{ij} x'_j \right) e_i$.
Puisque la décomposition de $u$ dans $\mathcal{B}$ est unique, on a $x_i = \sum_{j=1}^n P_{ij} x'_j$.
En notation matricielle, ceci s'écrit exactement : $X = P X'$.
En multipliant à gauche par $P^{-1}$, on obtient $X' = P^{-1} X$.

\textbf{4. Exemples et Cas Pathologiques :}
- *Cas trivial :* $\mathcal{B} = \mathcal{B}' \implies P = I_n$ (matrice identité).
- *Contre-exemple (famille liée) :* Si $\mathcal{B}'$ n'était pas libre, $P$ ne serait pas inversible et la relation $X = P X'$ ne pourrait pas être inversée, signifiant que les coordonnées dans $\mathcal{B}'$ ne seraient pas uniquement déterminées.


\section*{Exercice 05}
\subsection*{Énoncé}
Soit $E = \mathbb{R}_2[X]$ l'espace vectoriel des polynômes de degré inférieur ou égal à 2.
On considère la base canonique $\mathcal{B} = (1, X, X^2)$.
Soit $\mathcal{B}' = (P_0, P_1, P_2)$ la famille de polynômes définie par :
$P_0(X) = 1$
$P_1(X) = X + 1$
$P_2(X) = (X+1)^2 = X^2 + 2X + 1$

1. Démontrer que $\mathcal{B}'$ est une base de $E$.
2. Déterminer la matrice de passage $P$ de $\mathcal{B}$ à $\mathcal{B}'$.
3. Déterminer la matrice de passage inverse $P^{-1}$.
4. Soit l'endomorphisme $D$ de $E$ défini par $D(P) = P'$ (dérivation formelle).
   Écrire la matrice $A$ de $D$ dans la base canonique $\mathcal{B}$.
5. Écrire la matrice $A'$ de $D$ dans la base $\mathcal{B}'$ en utilisant la formule de changement de base.

\subsection*{Correction}
\textbf{1. Démontrer que $\mathcal{B}'$ est une base de $E$ :}
Les degrés des polynômes sont respectivement $\deg(P_0)=0$, $\deg(P_1)=1$, $\deg(P_2)=2$.
Une famille de polynômes non nuls de degrés échelonnés est toujours libre.
Comme la famille $(P_0, P_1, P_2)$ comporte 3 vecteurs et que $\dim(\mathbb{R}_2[X]) = 3$, c'est une base de $E$.

\textbf{2. Matrice de passage $P$ :}
On exprime les vecteurs de $\mathcal{B}'$ dans la base $\mathcal{B} = (1, X, X^2)$ :
$P_0(X) = 1 \cdot 1 + 0 \cdot X + 0 \cdot X^2 \implies \begin{pmatrix} 1 \\ 0 \\ 0 \end{pmatrix}$
$P_1(X) = 1 \cdot 1 + 1 \cdot X + 0 \cdot X^2 \implies \begin{pmatrix} 1 \\ 1 \\ 0 \end{pmatrix}$
$P_2(X) = 1 \cdot 1 + 2 \cdot X + 1 \cdot X^2 \implies \begin{pmatrix} 1 \\ 2 \\ 1 \end{pmatrix}$

La matrice de passage est constituée de ces vecteurs colonnes :
$P = \begin{pmatrix} 1 & 1 & 1 \\ 0 & 1 & 2 \\ 0 & 0 & 1 \end{pmatrix}$.

\textbf{3. Matrice inverse $P^{-1}$ :}
Pour inverser la matrice, on peut résoudre le système $Y = P X \iff X = P^{-1} Y$. Cela revient à exprimer les vecteurs de $\mathcal{B}$ en fonction de ceux de $\mathcal{B}'$.
$1 = P_0$
$X = P_1 - 1 = P_1 - P_0$
$X^2 = P_2 - 2X - 1 = P_2 - 2(P_1 - P_0) - P_0 = P_2 - 2P_1 + 2P_0 - P_0 = P_2 - 2P_1 + P_0$
On peut lire directement les coefficients sur les colonnes de $P^{-1}$ :
$P^{-1} = \begin{pmatrix} 1 & -1 & 1 \\ 0 & 1 & -2 \\ 0 & 0 & 1 \end{pmatrix}$.
*Vérification :* $\begin{pmatrix} 1 & 1 & 1 \\ 0 & 1 & 2 \\ 0 & 0 & 1 \end{pmatrix} \begin{pmatrix} 1 & -1 & 1 \\ 0 & 1 & -2 \\ 0 & 0 & 1 \end{pmatrix} = \begin{pmatrix} 1 & 0 & 0 \\ 0 & 1 & 0 \\ 0 & 0 & 1 \end{pmatrix} = I_3$.

\textbf{4. Matrice $A$ de $D$ dans $\mathcal{B}$ :}
$D(1) = 0 = 0 \cdot 1 + 0 \cdot X + 0 \cdot X^2$
$D(X) = 1 = 1 \cdot 1 + 0 \cdot X + 0 \cdot X^2$
$D(X^2) = 2X = 0 \cdot 1 + 2 \cdot X + 0 \cdot X^2$
Donc $A = \begin{pmatrix} 0 & 1 & 0 \\ 0 & 0 & 2 \\ 0 & 0 & 0 \end{pmatrix}$.

\textbf{5. Matrice $A'$ de $D$ dans $\mathcal{B}'$ :}
D'après la formule de changement de base pour un endomorphisme : $A' = P^{-1} A P$.
Calculons d'abord $AP$ :
$AP = \begin{pmatrix} 0 & 1 & 0 \\ 0 & 0 & 2 \\ 0 & 0 & 0 \end{pmatrix} \begin{pmatrix} 1 & 1 & 1 \\ 0 & 1 & 2 \\ 0 & 0 & 1 \end{pmatrix} = \begin{pmatrix} 0 & 1 & 2 \\ 0 & 0 & 2 \\ 0 & 0 & 0 \end{pmatrix}$.
Maintenant calculons $P^{-1}(AP)$ :
$A' = \begin{pmatrix} 1 & -1 & 1 \\ 0 & 1 & -2 \\ 0 & 0 & 1 \end{pmatrix} \begin{pmatrix} 0 & 1 & 2 \\ 0 & 0 & 2 \\ 0 & 0 & 0 \end{pmatrix} = \begin{pmatrix} 0 & 1 & 0 \\ 0 & 0 & 2 \\ 0 & 0 & 0 \end{pmatrix}$.
On remarque que $A' = A$. Ce résultat pouvait être prévu car $D(P_0) = 0$, $D(P_1) = 1 = P_0$, et $D(P_2) = 2(X+1) = 2P_1$.







\subsection*{Correction exhaustive (Protocole d'Exégèse)}

\textbf{1. Énoncé symbolique et Typage Chirurgical :}
Soit $E$ un $\mathbb{K}$-espace vectoriel de dimension finie $n$. Soient $\mathcal{B} = (e_1, \dots, e_n)$ et $\mathcal{B}' = (e'_1, \dots, e'_n)$ deux bases de $E$.
La matrice de passage $P = \operatorname{Mat}_{\mathcal{B}}(\mathcal{B}')$ a pour colonnes les composantes des vecteurs $e'_j$ exprimées dans la base $\mathcal{B}$.
Soit $u \in E$, on note $X = \operatorname{Mat}_{\mathcal{B}}(u)$ et $X' = \operatorname{Mat}_{\mathcal{B}'}(u)$.

\textbf{2. Démonstration de l'inversibilité de la matrice de passage (Zéro ellipse) :}
Puisque $\mathcal{B}'$ est une base, l'application linéaire $f : E \to E$ définie par $f(e_j) = e'_j$ pour tout $j \in \{1, \dots, n\}$ transforme une base en une base, c'est donc un automorphisme de $E$.
Or, par définition, $P = \operatorname{Mat}_{\mathcal{B}}(f)$.
Puisque $f$ est bijective, son déterminant est non nul : $\det(P) \neq 0$.
Ainsi, $P \in \operatorname{GL}_n(\mathbb{K})$ est inversible.

\textbf{3. Calcul de la relation de changement de base :}
On sait que $u = \sum_{j=1}^n x'_j e'_j$.
Remplaçons $e'_j$ par son expression dans $\mathcal{B}$ : $e'_j = \sum_{i=1}^n P_{ij} e_i$.
Alors $u = \sum_{j=1}^n x'_j \left( \sum_{i=1}^n P_{ij} e_i \right) = \sum_{i=1}^n \left( \sum_{j=1}^n P_{ij} x'_j \right) e_i$.
Puisque la décomposition de $u$ dans $\mathcal{B}$ est unique, on a $x_i = \sum_{j=1}^n P_{ij} x'_j$.
En notation matricielle, ceci s'écrit exactement : $X = P X'$.
En multipliant à gauche par $P^{-1}$, on obtient $X' = P^{-1} X$.

\textbf{4. Exemples et Cas Pathologiques :}
- *Cas trivial :* $\mathcal{B} = \mathcal{B}' \implies P = I_n$ (matrice identité).
- *Contre-exemple (famille liée) :* Si $\mathcal{B}'$ n'était pas libre, $P$ ne serait pas inversible et la relation $X = P X'$ ne pourrait pas être inversée, signifiant que les coordonnées dans $\mathcal{B}'$ ne seraient pas uniquement déterminées.


\section*{Exercice 06}
\subsection*{Énoncé}
Soit $E$ un espace vectoriel de dimension 3 rapporté à une base $\mathcal{B} = (e_1, e_2, e_3)$.
Soit $f$ l'endomorphisme de $E$ dont la matrice dans la base $\mathcal{B}$ est :
$A = \begin{pmatrix} 2 & -1 & -1 \\ -1 & 2 & -1 \\ -1 & -1 & 2 \end{pmatrix}$

1. On pose $u = e_1 + e_2 + e_3$. Calculer $f(u)$.
2. On pose $v = e_1 - e_2$ et $w = e_1 - e_3$. Calculer $f(v)$ et $f(w)$.
3. Montrer que $\mathcal{B}' = (u, v, w)$ est une base de $E$.
4. Écrire la matrice de passage $P$ de $\mathcal{B}$ à $\mathcal{B}'$.
5. Donner la matrice $A'$ de $f$ dans la base $\mathcal{B}'$ sans calculer $P^{-1}$.

\subsection*{Correction}
\textbf{1. Calcul de $f(u)$ :}
La colonne des coordonnées de $u$ dans $\mathcal{B}$ est $X = \begin{pmatrix} 1 \\ 1 \\ 1 \end{pmatrix}$.
La matrice colonne de $f(u)$ est :
$AX = \begin{pmatrix} 2 & -1 & -1 \\ -1 & 2 & -1 \\ -1 & -1 & 2 \end{pmatrix} \begin{pmatrix} 1 \\ 1 \\ 1 \end{pmatrix} = \begin{pmatrix} 2-1-1 \\ -1+2-1 \\ -1-1+2 \end{pmatrix} = \begin{pmatrix} 0 \\ 0 \\ 0 \end{pmatrix}$.
Donc $f(u) = 0_E$.

\textbf{2. Calcul de $f(v)$ et $f(w)$ :}
La colonne des coordonnées de $v$ dans $\mathcal{B}$ est $Y = \begin{pmatrix} 1 \\ -1 \\ 0 \end{pmatrix}$.
$AY = \begin{pmatrix} 2 & -1 & -1 \\ -1 & 2 & -1 \\ -1 & -1 & 2 \end{pmatrix} \begin{pmatrix} 1 \\ -1 \\ 0 \end{pmatrix} = \begin{pmatrix} 2-(-1)-0 \\ -1-2-0 \\ -1-(-1)-0 \end{pmatrix} = \begin{pmatrix} 3 \\ -3 \\ 0 \end{pmatrix} = 3 \begin{pmatrix} 1 \\ -1 \\ 0 \end{pmatrix}$.
Donc $f(v) = 3v$.

La colonne des coordonnées de $w$ dans $\mathcal{B}$ est $Z = \begin{pmatrix} 1 \\ 0 \\ -1 \end{pmatrix}$.
$AZ = \begin{pmatrix} 2 & -1 & -1 \\ -1 & 2 & -1 \\ -1 & -1 & 2 \end{pmatrix} \begin{pmatrix} 1 \\ 0 \\ -1 \end{pmatrix} = \begin{pmatrix} 2-0-(-1) \\ -1-0-(-1) \\ -1-0-2 \end{pmatrix} = \begin{pmatrix} 3 \\ 0 \\ -3 \end{pmatrix} = 3 \begin{pmatrix} 1 \\ 0 \\ -1 \end{pmatrix}$.
Donc $f(w) = 3w$.

\textbf{3. Montrer que $\mathcal{B}'$ est une base de $E$ :}
Il s'agit de montrer que la famille $(u, v, w)$ est libre. Soient $\lambda_1, \lambda_2, \lambda_3 \in \mathbb{R}$ tels que $\lambda_1 u + \lambda_2 v + \lambda_3 w = 0_E$.
On remplace $u, v, w$ par leurs expressions en fonction de $e_1, e_2, e_3$ :
$\lambda_1(e_1+e_2+e_3) + \lambda_2(e_1-e_2) + \lambda_3(e_1-e_3) = 0_E$
$(\lambda_1+\lambda_2+\lambda_3)e_1 + (\lambda_1-\lambda_2)e_2 + (\lambda_1-\lambda_3)e_3 = 0_E$.
Comme $(e_1, e_2, e_3)$ est une base, c'est une famille libre. On obtient le système :
$\begin{cases} \lambda_1 + \lambda_2 + \lambda_3 = 0 \\ \lambda_1 - \lambda_2 = 0 \\ \lambda_1 - \lambda_3 = 0 \end{cases}$
De la 2e et 3e équation, on tire $\lambda_2 = \lambda_1$ et $\lambda_3 = \lambda_1$.
En remplaçant dans la 1ère : $\lambda_1 + \lambda_1 + \lambda_1 = 0 \implies 3\lambda_1 = 0 \implies \lambda_1 = 0$.
D'où $\lambda_2 = 0$ et $\lambda_3 = 0$.
La famille $(u, v, w)$ est libre. Étant de cardinal 3 dans un espace de dimension 3, c'est une base.

\textbf{4. Matrice de passage $P$ :}
Les colonnes de $P$ sont les coordonnées de $u, v, w$ exprimées dans la base $(e_1, e_2, e_3)$ :
$P = \begin{pmatrix} 1 & 1 & 1 \\ 1 & -1 & 0 \\ 1 & 0 & -1 \end{pmatrix}$.

\textbf{5. Matrice $A'$ de $f$ dans $\mathcal{B}'$ :}
La matrice $A'$ représente l'endomorphisme $f$ dans la base $\mathcal{B}' = (u, v, w)$.
Ses colonnes sont les coordonnées de $f(u), f(v), f(w)$ exprimées dans la base $(u, v, w)$.
D'après les questions 1 et 2, on a :
$f(u) = 0_E = 0\cdot u + 0\cdot v + 0\cdot w$
$f(v) = 3v = 0\cdot u + 3\cdot v + 0\cdot w$
$f(w) = 3w = 0\cdot u + 0\cdot v + 3\cdot w$
Ainsi, la matrice s'écrit directement :
$A' = \begin{pmatrix} 0 & 0 & 0 \\ 0 & 3 & 0 \\ 0 & 0 & 3 \end{pmatrix}$.
Cette méthode est bien plus rapide que de calculer explicitement $P^{-1} A P$.







\subsection*{Correction exhaustive (Protocole d'Exégèse)}

\textbf{1. Énoncé symbolique et Typage Chirurgical :}
Soit $E$ un $\mathbb{K}$-espace vectoriel de dimension finie $n$. Soient $\mathcal{B} = (e_1, \dots, e_n)$ et $\mathcal{B}' = (e'_1, \dots, e'_n)$ deux bases de $E$.
La matrice de passage $P = \operatorname{Mat}_{\mathcal{B}}(\mathcal{B}')$ a pour colonnes les composantes des vecteurs $e'_j$ exprimées dans la base $\mathcal{B}$.
Soit $u \in E$, on note $X = \operatorname{Mat}_{\mathcal{B}}(u)$ et $X' = \operatorname{Mat}_{\mathcal{B}'}(u)$.

\textbf{2. Démonstration de l'inversibilité de la matrice de passage (Zéro ellipse) :}
Puisque $\mathcal{B}'$ est une base, l'application linéaire $f : E \to E$ définie par $f(e_j) = e'_j$ pour tout $j \in \{1, \dots, n\}$ transforme une base en une base, c'est donc un automorphisme de $E$.
Or, par définition, $P = \operatorname{Mat}_{\mathcal{B}}(f)$.
Puisque $f$ est bijective, son déterminant est non nul : $\det(P) \neq 0$.
Ainsi, $P \in \operatorname{GL}_n(\mathbb{K})$ est inversible.

\textbf{3. Calcul de la relation de changement de base :}
On sait que $u = \sum_{j=1}^n x'_j e'_j$.
Remplaçons $e'_j$ par son expression dans $\mathcal{B}$ : $e'_j = \sum_{i=1}^n P_{ij} e_i$.
Alors $u = \sum_{j=1}^n x'_j \left( \sum_{i=1}^n P_{ij} e_i \right) = \sum_{i=1}^n \left( \sum_{j=1}^n P_{ij} x'_j \right) e_i$.
Puisque la décomposition de $u$ dans $\mathcal{B}$ est unique, on a $x_i = \sum_{j=1}^n P_{ij} x'_j$.
En notation matricielle, ceci s'écrit exactement : $X = P X'$.
En multipliant à gauche par $P^{-1}$, on obtient $X' = P^{-1} X$.

\textbf{4. Exemples et Cas Pathologiques :}
- *Cas trivial :* $\mathcal{B} = \mathcal{B}' \implies P = I_n$ (matrice identité).
- *Contre-exemple (famille liée) :* Si $\mathcal{B}'$ n'était pas libre, $P$ ne serait pas inversible et la relation $X = P X'$ ne pourrait pas être inversée, signifiant que les coordonnées dans $\mathcal{B}'$ ne seraient pas uniquement déterminées.


\section*{Exercice 07}
\subsection*{Énoncé}
Soit $M = \begin{pmatrix} A & B \\ 0 & C \end{pmatrix}$ une matrice par blocs, où $A \in \mathcal{M}_n(\mathbb{R})$, $C \in \mathcal{M}_p(\mathbb{R})$ et $B \in \mathcal{M}_{n,p}(\mathbb{R})$.
1. Calculer $M^2$ puis $M^3$ sous forme de matrices par blocs.
2. Démontrer par récurrence sur $k \in \mathbb{N}^*$ que $M^k = \begin{pmatrix} A^k & X_k \\ 0 & C^k \end{pmatrix}$ où on précisera une relation de récurrence pour la matrice $X_k$.
3. On pose $A = \begin{pmatrix} 2 & 0 \\ 0 & 2 \end{pmatrix}$, $C = \begin{pmatrix} 3 \end{pmatrix}$ et $B = \begin{pmatrix} 1 \\ 1 \end{pmatrix}$. Calculer explicitement la matrice $M^k$ pour tout $k \ge 1$.

\subsection*{Correction}
\textbf{1. Calcul de $M^2$ et $M^3$ :}
Les produits matriciels s'effectuent sur les blocs de la même manière que sur des scalaires (à condition que les dimensions soient compatibles, ce qui est le cas ici).
$M^2 = \begin{pmatrix} A & B \\ 0 & C \end{pmatrix} \begin{pmatrix} A & B \\ 0 & C \end{pmatrix} = \begin{pmatrix} A \cdot A + B \cdot 0 & A \cdot B + B \cdot C \\ 0 \cdot A + C \cdot 0 & 0 \cdot B + C \cdot C \end{pmatrix} = \begin{pmatrix} A^2 & AB+BC \\ 0 & C^2 \end{pmatrix}$.

Pour $M^3$ :
$M^3 = M^2 M = \begin{pmatrix} A^2 & AB+BC \\ 0 & C^2 \end{pmatrix} \begin{pmatrix} A & B \\ 0 & C \end{pmatrix} = \begin{pmatrix} A^2 \cdot A + (AB+BC) \cdot 0 & A^2 \cdot B + (AB+BC)C \\ 0 \cdot A + C^2 \cdot 0 & 0 \cdot B + C^2 \cdot C \end{pmatrix} = \begin{pmatrix} A^3 & A^2B + ABC + BC^2 \\ 0 & C^3 \end{pmatrix}$.

\textbf{2. Démonstration par récurrence :}
Soit la propriété $P(k)$ : "$M^k = \begin{pmatrix} A^k & X_k \\ 0 & C^k \end{pmatrix}$".
*Initialisation :* Pour $k=1$, $M^1 = M = \begin{pmatrix} A & B \\ 0 & C \end{pmatrix}$, ce qui correspond bien à la forme avec $X_1 = B$.
*Hérédité :* Supposons la propriété vraie pour un certain entier $k \ge 1$.
Montrons-la au rang $k+1$.
$M^{k+1} = M^k M = \begin{pmatrix} A^k & X_k \\ 0 & C^k \end{pmatrix} \begin{pmatrix} A & B \\ 0 & C \end{pmatrix} = \begin{pmatrix} A^k \cdot A + X_k \cdot 0 & A^k \cdot B + X_k \cdot C \\ 0 \cdot A + C^k \cdot 0 & 0 \cdot B + C^k \cdot C \end{pmatrix} = \begin{pmatrix} A^{k+1} & A^kB + X_k C \\ 0 & C^{k+1} \end{pmatrix}$.
La propriété est donc vraie au rang $k+1$ en posant $X_{k+1} = A^k B + X_k C$.
*Conclusion :* Par le principe de récurrence, pour tout $k \ge 1$, $M^k = \begin{pmatrix} A^k & X_k \\ 0 & C^k \end{pmatrix}$.

\textbf{3. Calcul de $M^k$ pour un cas spécifique :}
On a $A = 2I_2$, $C = (3)$, et $B = \begin{pmatrix} 1 \\ 1 \end{pmatrix}$.
Il est immédiat que $A^k = 2^k I_2$ et $C^k = (3^k)$.
Il reste à calculer $X_k$. Trouvons une formule explicite pour $X_k$.
D'après la relation de récurrence obtenue en question 2, on peut aussi l'écrire de l'autre côté :
$M^{k+1} = M M^k \implies X_{k+1} = A X_k + B C^k$.
Puisque $A = 2I_2$, on a : $X_{k+1} = 2 X_k + B 3^k$.
Calculons les premiers termes :
$X_1 = B$
$X_2 = 2B + 3B = 5B$
$X_3 = 2(5B) + 3^2 B = 10B + 9B = 19B$.
On peut résoudre cette suite arithmético-géométrique matricielle. Soit $X_k = x_k B$.
La suite scalaire $(x_k)$ vérifie $x_{k+1} = 2x_k + 3^k$ avec $x_1 = 1$.
Cherchons une solution particulière sous la forme $x_k = \lambda 3^k$.
$\lambda 3^{k+1} = 2 \lambda 3^k + 3^k \implies 3\lambda = 2\lambda + 1 \implies \lambda = 1$.
Donc $x_k = 3^k + \alpha 2^k$.
Pour $k=1$, $x_1 = 1 = 3^1 + \alpha 2^1 \implies 1 = 3 + 2\alpha \implies 2\alpha = -2 \implies \alpha = -1$.
Donc $x_k = 3^k - 2^k$.
Ainsi, $X_k = (3^k - 2^k) B = (3^k - 2^k) \begin{pmatrix} 1 \\ 1 \end{pmatrix} = \begin{pmatrix} 3^k - 2^k \\ 3^k - 2^k \end{pmatrix}$.
Finalement, pour tout $k \ge 1$ :
$M^k = \begin{pmatrix} 2^k & 0 & 3^k - 2^k \\ 0 & 2^k & 3^k - 2^k \\ 0 & 0 & 3^k \end{pmatrix}$.







\subsection*{Correction exhaustive (Protocole d'Exégèse)}

\textbf{1. Énoncé symbolique et Typage Chirurgical :}
Soit $E$ un $\mathbb{K}$-espace vectoriel de dimension finie $n$. Soient $\mathcal{B} = (e_1, \dots, e_n)$ et $\mathcal{B}' = (e'_1, \dots, e'_n)$ deux bases de $E$.
La matrice de passage $P = \operatorname{Mat}_{\mathcal{B}}(\mathcal{B}')$ a pour colonnes les composantes des vecteurs $e'_j$ exprimées dans la base $\mathcal{B}$.
Soit $u \in E$, on note $X = \operatorname{Mat}_{\mathcal{B}}(u)$ et $X' = \operatorname{Mat}_{\mathcal{B}'}(u)$.

\textbf{2. Démonstration de l'inversibilité de la matrice de passage (Zéro ellipse) :}
Puisque $\mathcal{B}'$ est une base, l'application linéaire $f : E \to E$ définie par $f(e_j) = e'_j$ pour tout $j \in \{1, \dots, n\}$ transforme une base en une base, c'est donc un automorphisme de $E$.
Or, par définition, $P = \operatorname{Mat}_{\mathcal{B}}(f)$.
Puisque $f$ est bijective, son déterminant est non nul : $\det(P) \neq 0$.
Ainsi, $P \in \operatorname{GL}_n(\mathbb{K})$ est inversible.

\textbf{3. Calcul de la relation de changement de base :}
On sait que $u = \sum_{j=1}^n x'_j e'_j$.
Remplaçons $e'_j$ par son expression dans $\mathcal{B}$ : $e'_j = \sum_{i=1}^n P_{ij} e_i$.
Alors $u = \sum_{j=1}^n x'_j \left( \sum_{i=1}^n P_{ij} e_i \right) = \sum_{i=1}^n \left( \sum_{j=1}^n P_{ij} x'_j \right) e_i$.
Puisque la décomposition de $u$ dans $\mathcal{B}$ est unique, on a $x_i = \sum_{j=1}^n P_{ij} x'_j$.
En notation matricielle, ceci s'écrit exactement : $X = P X'$.
En multipliant à gauche par $P^{-1}$, on obtient $X' = P^{-1} X$.

\textbf{4. Exemples et Cas Pathologiques :}
- *Cas trivial :* $\mathcal{B} = \mathcal{B}' \implies P = I_n$ (matrice identité).
- *Contre-exemple (famille liée) :* Si $\mathcal{B}'$ n'était pas libre, $P$ ne serait pas inversible et la relation $X = P X'$ ne pourrait pas être inversée, signifiant que les coordonnées dans $\mathcal{B}'$ ne seraient pas uniquement déterminées.


\section*{Exercice 08}
\subsection*{Énoncé}
Soient $E_1, E_2$ deux espaces vectoriels de dimensions finies $n$ et $p$. On considère l'espace vectoriel produit $E = E_1 \times E_2$.
Soit $f_1 \in \mathcal{L}(E_1)$ et $f_2 \in \mathcal{L}(E_2)$.
On définit l'application $f : E \to E$ par $f(x,y) = (f_1(x), f_2(y))$.
Soit $\mathcal{B}_1$ une base de $E_1$ et $\mathcal{B}_2$ une base de $E_2$.
1. Montrer que $f$ est un endomorphisme de $E$.
2. Construire une base $\mathcal{B}$ de $E$ à partir de $\mathcal{B}_1$ et $\mathcal{B}_2$. Quelle est la dimension de $E$ ?
3. Montrer que la matrice $M$ de $f$ dans la base $\mathcal{B}$ est une matrice diagonale par blocs.
4. En déduire une expression du polynôme caractéristique de $M$ en fonction de ceux des blocs. *(Bien que le polynôme caractéristique soit formellement abordé au Jalon 29, il s'agit ici simplement d'appliquer la définition par le déterminant d'une matrice par blocs).*

\subsection*{Correction}
\textbf{1. Endomorphisme :}
$E_1 \times E_2$ est un espace vectoriel. $f$ va de $E$ dans $E$.
Montrons la linéarité. Soient $(x,y) \in E$, $(x',y') \in E$ et $\lambda \in \mathbb{R}$.
$f((x,y) + \lambda(x',y')) = f(x+\lambda x', y+\lambda y') = (f_1(x+\lambda x'), f_2(y+\lambda y'))$.
Par linéarité de $f_1$ et $f_2$ :
$= (f_1(x) + \lambda f_1(x'), f_2(y) + \lambda f_2(y')) = (f_1(x), f_2(y)) + \lambda(f_1(x'), f_2(y')) = f(x,y) + \lambda f(x',y')$.
$f$ est donc linéaire, c'est un endomorphisme.

\textbf{2. Base de l'espace produit :}
Soit $\mathcal{B}_1 = (e_1, ..., e_n)$ et $\mathcal{B}_2 = (u_1, ..., u_p)$.
On forme la famille $\mathcal{B} = ((e_1, 0), ..., (e_n, 0), (0, u_1), ..., (0, u_p))$.
Montrons qu'elle est génératrice : pour tout $(x,y) \in E$, on peut décomposer $x = \sum x_i e_i$ et $y = \sum y_j u_j$.
Alors $(x,y) = (x,0) + (0,y) = \sum x_i (e_i, 0) + \sum y_j (0, u_j)$.
Montrons qu'elle est libre : $\sum \lambda_i (e_i, 0) + \sum \mu_j (0, u_j) = (0,0) \implies (\sum \lambda_i e_i, \sum \mu_j u_j) = (0,0)$.
Ceci implique $\sum \lambda_i e_i = 0$ et $\sum \mu_j u_j = 0$. Comme $\mathcal{B}_1$ et $\mathcal{B}_2$ sont libres, tous les $\lambda_i$ et $\mu_j$ sont nuls.
Donc $\mathcal{B}$ est une base de $E$, et $\dim(E) = n + p$.

\textbf{3. Matrice de $f$ dans $\mathcal{B}$ :}
Évaluons $f$ sur les vecteurs de base :
$f(e_i, 0) = (f_1(e_i), f_2(0)) = (f_1(e_i), 0)$.
Comme $f_1(e_i) \in E_1$, ses coordonnées ne s'expriment que sur la première partie de la base $\mathcal{B}$ (les $(e_k, 0)$), et les coefficients sur les $(0, u_j)$ sont nuls.
De même, $f(0, u_j) = (f_1(0), f_2(u_j)) = (0, f_2(u_j))$, dont les coordonnées ne s'expriment que sur la seconde partie de la base.
Si l'on pose $A = \text{Mat}_{\mathcal{B}_1}(f_1)$ et $B = \text{Mat}_{\mathcal{B}_2}(f_2)$, la matrice $M$ s'écrit formellement par blocs :
$M = \begin{pmatrix} A & 0 \\ 0 & B \end{pmatrix}$.

\textbf{4. Déterminant et polynôme :}
Le déterminant d'une matrice diagonale par blocs est le produit des déterminants des blocs diagonaux.
Le polynôme caractéristique est $P_M(\lambda) = \det(M - \lambda I_{n+p})$.
Or, $M - \lambda I_{n+p} = \begin{pmatrix} A - \lambda I_n & 0 \\ 0 & B - \lambda I_p \end{pmatrix}$.
Donc $P_M(\lambda) = \det(A - \lambda I_n) \times \det(B - \lambda I_p) = P_A(\lambda) \times P_B(\lambda)$.







\subsection*{Correction exhaustive (Protocole d'Exégèse)}

\textbf{1. Énoncé symbolique et Typage Chirurgical :}
Soit $E$ un $\mathbb{K}$-espace vectoriel de dimension finie $n$. Soient $\mathcal{B} = (e_1, \dots, e_n)$ et $\mathcal{B}' = (e'_1, \dots, e'_n)$ deux bases de $E$.
La matrice de passage $P = \operatorname{Mat}_{\mathcal{B}}(\mathcal{B}')$ a pour colonnes les composantes des vecteurs $e'_j$ exprimées dans la base $\mathcal{B}$.
Soit $u \in E$, on note $X = \operatorname{Mat}_{\mathcal{B}}(u)$ et $X' = \operatorname{Mat}_{\mathcal{B}'}(u)$.

\textbf{2. Démonstration de l'inversibilité de la matrice de passage (Zéro ellipse) :}
Puisque $\mathcal{B}'$ est une base, l'application linéaire $f : E \to E$ définie par $f(e_j) = e'_j$ pour tout $j \in \{1, \dots, n\}$ transforme une base en une base, c'est donc un automorphisme de $E$.
Or, par définition, $P = \operatorname{Mat}_{\mathcal{B}}(f)$.
Puisque $f$ est bijective, son déterminant est non nul : $\det(P) \neq 0$.
Ainsi, $P \in \operatorname{GL}_n(\mathbb{K})$ est inversible.

\textbf{3. Calcul de la relation de changement de base :}
On sait que $u = \sum_{j=1}^n x'_j e'_j$.
Remplaçons $e'_j$ par son expression dans $\mathcal{B}$ : $e'_j = \sum_{i=1}^n P_{ij} e_i$.
Alors $u = \sum_{j=1}^n x'_j \left( \sum_{i=1}^n P_{ij} e_i \right) = \sum_{i=1}^n \left( \sum_{j=1}^n P_{ij} x'_j \right) e_i$.
Puisque la décomposition de $u$ dans $\mathcal{B}$ est unique, on a $x_i = \sum_{j=1}^n P_{ij} x'_j$.
En notation matricielle, ceci s'écrit exactement : $X = P X'$.
En multipliant à gauche par $P^{-1}$, on obtient $X' = P^{-1} X$.

\textbf{4. Exemples et Cas Pathologiques :}
- *Cas trivial :* $\mathcal{B} = \mathcal{B}' \implies P = I_n$ (matrice identité).
- *Contre-exemple (famille liée) :* Si $\mathcal{B}'$ n'était pas libre, $P$ ne serait pas inversible et la relation $X = P X'$ ne pourrait pas être inversée, signifiant que les coordonnées dans $\mathcal{B}'$ ne seraient pas uniquement déterminées.


\section*{Exercice 09}
\subsection*{Énoncé}
Soient $A, B \in \mathcal{M}_n(\mathbb{R})$.
On considère la matrice par blocs de taille $2n \times 2n$ : $M = \begin{pmatrix} A & B \\ B & A \end{pmatrix}$.
1. En introduisant la matrice de passage par blocs $P = \begin{pmatrix} I_n & I_n \\ I_n & -I_n \end{pmatrix}$, calculer l'inverse $P^{-1}$.
2. Calculer le produit $P^{-1} M P$.
3. En déduire une condition nécessaire et suffisante pour que $M$ soit inversible, exprimée à l'aide de $A+B$ et $A-B$.
4. Si cette condition est remplie, exprimer $M^{-1}$ sous forme de matrice par blocs en fonction de $(A+B)^{-1}$ et $(A-B)^{-1}$.

\subsection*{Correction}
\textbf{1. Inversion de $P$ :}
Cherchons $P^{-1} = \begin{pmatrix} X & Y \\ Z & W \end{pmatrix}$ telle que $P P^{-1} = \begin{pmatrix} I_n & 0 \\ 0 & I_n \end{pmatrix}$.
$\begin{pmatrix} I_n & I_n \\ I_n & -I_n \end{pmatrix} \begin{pmatrix} X & Y \\ Z & W \end{pmatrix} = \begin{pmatrix} X+Z & Y+W \\ X-Z & Y-W \end{pmatrix} = \begin{pmatrix} I_n & 0 \\ 0 & I_n \end{pmatrix}$.
On résout les systèmes :
$X+Z = I_n$ et $X-Z = 0 \implies X=Z$. Donc $2X = I_n \implies X = Z = \frac{1}{2}I_n$.
$Y+W = 0 \implies W = -Y$. $Y-W = I_n \implies 2Y = I_n \implies Y = \frac{1}{2}I_n$, $W = -\frac{1}{2}I_n$.
Donc $P^{-1} = \frac{1}{2} \begin{pmatrix} I_n & I_n \\ I_n & -I_n \end{pmatrix} = \frac{1}{2} P$.
(On remarque que $P^2 = 2 I_{2n}$).

\textbf{2. Calcul de $P^{-1} M P$ :}
$MP = \begin{pmatrix} A & B \\ B & A \end{pmatrix} \begin{pmatrix} I_n & I_n \\ I_n & -I_n \end{pmatrix} = \begin{pmatrix} A+B & A-B \\ B+A & B-A \end{pmatrix}$.
$P^{-1} M P = \frac{1}{2} \begin{pmatrix} I_n & I_n \\ I_n & -I_n \end{pmatrix} \begin{pmatrix} A+B & A-B \\ A+B & -(A-B) \end{pmatrix} = \frac{1}{2} \begin{pmatrix} (A+B)+(A+B) & (A-B)-(A-B) \\ (A+B)-(A+B) & (A-B)-(-(A-B)) \end{pmatrix} = \frac{1}{2} \begin{pmatrix} 2(A+B) & 0 \\ 0 & 2(A-B) \end{pmatrix} = \begin{pmatrix} A+B & 0 \\ 0 & A-B \end{pmatrix}$.

\textbf{3. Condition d'inversibilité :}
On a obtenu $P^{-1} M P = \begin{pmatrix} A+B & 0 \\ 0 & A-B \end{pmatrix}$.
Soit $D$ cette matrice diagonale par blocs. Comme $P$ est inversible, $M$ est inversible si et seulement si $D$ l'est.
Or, l'inversibilité d'une matrice diagonale par blocs équivaut à l'inversibilité de chacun de ses blocs diagonaux.
Donc $M$ est inversible $\iff (A+B)$ est inversible ET $(A-B)$ est inversible.

\textbf{4. Calcul de $M^{-1}$ :}
Si la condition est remplie, $D^{-1} = \begin{pmatrix} (A+B)^{-1} & 0 \\ 0 & (A-B)^{-1} \end{pmatrix}$.
Puisque $D = P^{-1} M P$, on a $M = P D P^{-1}$, d'où $M^{-1} = P D^{-1} P^{-1}$.
$M^{-1} = \begin{pmatrix} I_n & I_n \\ I_n & -I_n \end{pmatrix} \begin{pmatrix} (A+B)^{-1} & 0 \\ 0 & (A-B)^{-1} \end{pmatrix} \left( \frac{1}{2} \begin{pmatrix} I_n & I_n \\ I_n & -I_n \end{pmatrix} \right)$
$M^{-1} = \frac{1}{2} \begin{pmatrix} (A+B)^{-1} & (A-B)^{-1} \\ (A+B)^{-1} & -(A-B)^{-1} \end{pmatrix} \begin{pmatrix} I_n & I_n \\ I_n & -I_n \end{pmatrix} = \frac{1}{2} \begin{pmatrix} (A+B)^{-1} + (A-B)^{-1} & (A+B)^{-1} - (A-B)^{-1} \\ (A+B)^{-1} - (A-B)^{-1} & (A+B)^{-1} + (A-B)^{-1} \end{pmatrix}$.
On pose $S = \frac{1}{2}((A+B)^{-1} + (A-B)^{-1})$ et $T = \frac{1}{2}((A+B)^{-1} - (A-B)^{-1})$.
Alors $M^{-1} = \begin{pmatrix} S & T \\ T & S \end{pmatrix}$. L'inverse possède la même structure de blocs que $M$.







\subsection*{Correction exhaustive (Protocole d'Exégèse)}

\textbf{1. Énoncé symbolique et Typage Chirurgical :}
Soit $E$ un $\mathbb{K}$-espace vectoriel de dimension finie $n$. Soient $\mathcal{B} = (e_1, \dots, e_n)$ et $\mathcal{B}' = (e'_1, \dots, e'_n)$ deux bases de $E$.
La matrice de passage $P = \operatorname{Mat}_{\mathcal{B}}(\mathcal{B}')$ a pour colonnes les composantes des vecteurs $e'_j$ exprimées dans la base $\mathcal{B}$.
Soit $u \in E$, on note $X = \operatorname{Mat}_{\mathcal{B}}(u)$ et $X' = \operatorname{Mat}_{\mathcal{B}'}(u)$.

\textbf{2. Démonstration de l'inversibilité de la matrice de passage (Zéro ellipse) :}
Puisque $\mathcal{B}'$ est une base, l'application linéaire $f : E \to E$ définie par $f(e_j) = e'_j$ pour tout $j \in \{1, \dots, n\}$ transforme une base en une base, c'est donc un automorphisme de $E$.
Or, par définition, $P = \operatorname{Mat}_{\mathcal{B}}(f)$.
Puisque $f$ est bijective, son déterminant est non nul : $\det(P) \neq 0$.
Ainsi, $P \in \operatorname{GL}_n(\mathbb{K})$ est inversible.

\textbf{3. Calcul de la relation de changement de base :}
On sait que $u = \sum_{j=1}^n x'_j e'_j$.
Remplaçons $e'_j$ par son expression dans $\mathcal{B}$ : $e'_j = \sum_{i=1}^n P_{ij} e_i$.
Alors $u = \sum_{j=1}^n x'_j \left( \sum_{i=1}^n P_{ij} e_i \right) = \sum_{i=1}^n \left( \sum_{j=1}^n P_{ij} x'_j \right) e_i$.
Puisque la décomposition de $u$ dans $\mathcal{B}$ est unique, on a $x_i = \sum_{j=1}^n P_{ij} x'_j$.
En notation matricielle, ceci s'écrit exactement : $X = P X'$.
En multipliant à gauche par $P^{-1}$, on obtient $X' = P^{-1} X$.

\textbf{4. Exemples et Cas Pathologiques :}
- *Cas trivial :* $\mathcal{B} = \mathcal{B}' \implies P = I_n$ (matrice identité).
- *Contre-exemple (famille liée) :* Si $\mathcal{B}'$ n'était pas libre, $P$ ne serait pas inversible et la relation $X = P X'$ ne pourrait pas être inversée, signifiant que les coordonnées dans $\mathcal{B}'$ ne seraient pas uniquement déterminées.


\section*{Exercice 10}
\subsection*{Énoncé}
Soit $E$ un espace vectoriel de dimension $n$ et $f, g \in \mathcal{L}(E)$ tels que $E = \ker(f) \oplus \text{Im}(f)$.
1. Montrer qu'il existe une base $\mathcal{B}$ de $E$ telle que la matrice $A$ de $f$ dans $\mathcal{B}$ soit de la forme $A = \begin{pmatrix} A_1 & 0 \\ 0 & 0 \end{pmatrix}$, où $A_1$ est une matrice carrée inversible de taille $r = \text{rg}(f)$.
2. On suppose de plus que $f \circ g = g \circ f = f$. Soit $B = \begin{pmatrix} B_1 & B_2 \\ B_3 & B_4 \end{pmatrix}$ la matrice de $g$ dans cette même base $\mathcal{B}$ (avec les mêmes tailles de blocs).
Traduire la condition $f \circ g = g \circ f = f$ sur les blocs de $B$.
3. En déduire que $\ker(g) \subset \ker(f)$.

\subsection*{Correction}
\textbf{1. Forme bloc pour la matrice de $f$ :}
Par hypothèse, $E = \text{Im}(f) \oplus \ker(f)$.
Soit $r = \dim(\text{Im}(f))$. D'après le théorème du rang, $\dim(\ker(f)) = n - r$.
Soit $(e_1, ..., e_r)$ une base de $\text{Im}(f)$ et $(e_{r+1}, ..., e_n)$ une base de $\ker(f)$.
Puisque la somme est directe, $\mathcal{B} = (e_1, ..., e_r, e_{r+1}, ..., e_n)$ est une base de $E$.
Écrivons la matrice $A$ de $f$ dans $\mathcal{B}$.
Pour $j > r$, $e_j \in \ker(f)$, donc $f(e_j) = 0$. Les $n-r$ dernières colonnes de $A$ sont nulles.
Pour $j \le r$, $f(e_j) \in \text{Im}(f)$, car $\text{Im}(f)$ est stable par $f$. Donc $f(e_j)$ s'exprime uniquement en fonction de $(e_1, ..., e_r)$. Les coefficients correspondants à $(e_{r+1}, ..., e_n)$ sont nuls.
Ainsi, $A = \begin{pmatrix} A_1 & 0 \\ 0 & 0 \end{pmatrix}$, où $A_1 \in \mathcal{M}_r(\mathbb{R})$.
Montrons que $A_1$ est inversible. Soit $u \in \text{Im}(f)$ dont le vecteur colonne dans $(e_1, ..., e_r)$ est $X$.
Si $A_1 X = 0$, cela signifie que $f(u) = 0$, donc $u \in \ker(f)$. Or $u \in \text{Im}(f)$. Puisque la somme est directe, $u = 0$.
Le noyau de l'endomorphisme induit par $A_1$ est réduit à $0$, donc $A_1$ est inversible.

\textbf{2. Condition de commutation avec $g$ :}
On a $A = \begin{pmatrix} A_1 & 0 \\ 0 & 0 \end{pmatrix}$ et $B = \begin{pmatrix} B_1 & B_2 \\ B_3 & B_4 \end{pmatrix}$.
L'équation matricielle est $AB = BA = A$.
Calculons les produits par blocs :
$AB = \begin{pmatrix} A_1 B_1 & A_1 B_2 \\ 0 & 0 \end{pmatrix}$
$BA = \begin{pmatrix} B_1 A_1 & 0 \\ B_3 A_1 & 0 \end{pmatrix}$
On identifie les blocs avec $A$ :
- $A_1 B_1 = A_1 \implies A_1(B_1 - I_r) = 0$. Comme $A_1$ est inversible, on multiplie par $A_1^{-1}$ à gauche : $B_1 = I_r$.
- $A_1 B_2 = 0 \implies B_2 = 0$ (pour la même raison).
- $B_1 A_1 = A_1 \implies B_1 = I_r$ (cohérent).
- $B_3 A_1 = 0 \implies B_3 A_1 A_1^{-1} = 0 \implies B_3 = 0$.
Le bloc $B_4$ n'est soumis à aucune contrainte par ces équations.
Donc la matrice $B$ est de la forme $B = \begin{pmatrix} I_r & 0 \\ 0 & B_4 \end{pmatrix}$.

\textbf{3. Déduction sur le noyau :}
Soit $x \in \ker(g)$. Son vecteur coordonnée $X = \begin{pmatrix} X_1 \\ X_2 \end{pmatrix}$ vérifie $BX = 0$.
$\begin{pmatrix} I_r & 0 \\ 0 & B_4 \end{pmatrix} \begin{pmatrix} X_1 \\ X_2 \end{pmatrix} = \begin{pmatrix} X_1 \\ B_4 X_2 \end{pmatrix} = \begin{pmatrix} 0 \\ 0 \end{pmatrix}$.
On en déduit que $X_1 = 0$.
Calculons l'image de $x$ par $f$ : le vecteur coordonnée est $AX$.
$AX = \begin{pmatrix} A_1 & 0 \\ 0 & 0 \end{pmatrix} \begin{pmatrix} 0 \\ X_2 \end{pmatrix} = \begin{pmatrix} A_1 \cdot 0 + 0 \cdot X_2 \\ 0 \cdot 0 + 0 \cdot X_2 \end{pmatrix} = \begin{pmatrix} 0 \\ 0 \end{pmatrix}$.
Donc $f(x) = 0$, c'est-à-dire $x \in \ker(f)$.
On a bien $\ker(g) \subset \ker(f)$.







\subsection*{Correction exhaustive (Protocole d'Exégèse)}

\textbf{1. Énoncé symbolique et Typage Chirurgical :}
Soit $E$ un $\mathbb{K}$-espace vectoriel de dimension finie $n$. Soient $\mathcal{B} = (e_1, \dots, e_n)$ et $\mathcal{B}' = (e'_1, \dots, e'_n)$ deux bases de $E$.
La matrice de passage $P = \operatorname{Mat}_{\mathcal{B}}(\mathcal{B}')$ a pour colonnes les composantes des vecteurs $e'_j$ exprimées dans la base $\mathcal{B}$.
Soit $u \in E$, on note $X = \operatorname{Mat}_{\mathcal{B}}(u)$ et $X' = \operatorname{Mat}_{\mathcal{B}'}(u)$.

\textbf{2. Démonstration de l'inversibilité de la matrice de passage (Zéro ellipse) :}
Puisque $\mathcal{B}'$ est une base, l'application linéaire $f : E \to E$ définie par $f(e_j) = e'_j$ pour tout $j \in \{1, \dots, n\}$ transforme une base en une base, c'est donc un automorphisme de $E$.
Or, par définition, $P = \operatorname{Mat}_{\mathcal{B}}(f)$.
Puisque $f$ est bijective, son déterminant est non nul : $\det(P) \neq 0$.
Ainsi, $P \in \operatorname{GL}_n(\mathbb{K})$ est inversible.

\textbf{3. Calcul de la relation de changement de base :}
On sait que $u = \sum_{j=1}^n x'_j e'_j$.
Remplaçons $e'_j$ par son expression dans $\mathcal{B}$ : $e'_j = \sum_{i=1}^n P_{ij} e_i$.
Alors $u = \sum_{j=1}^n x'_j \left( \sum_{i=1}^n P_{ij} e_i \right) = \sum_{i=1}^n \left( \sum_{j=1}^n P_{ij} x'_j \right) e_i$.
Puisque la décomposition de $u$ dans $\mathcal{B}$ est unique, on a $x_i = \sum_{j=1}^n P_{ij} x'_j$.
En notation matricielle, ceci s'écrit exactement : $X = P X'$.
En multipliant à gauche par $P^{-1}$, on obtient $X' = P^{-1} X$.

\textbf{4. Exemples et Cas Pathologiques :}
- *Cas trivial :* $\mathcal{B} = \mathcal{B}' \implies P = I_n$ (matrice identité).
- *Contre-exemple (famille liée) :* Si $\mathcal{B}'$ n'était pas libre, $P$ ne serait pas inversible et la relation $X = P X'$ ne pourrait pas être inversée, signifiant que les coordonnées dans $\mathcal{B}'$ ne seraient pas uniquement déterminées.


\chapter{Travaux Pratiques}
\section*{TP 01}
\subsection*{Objectif}
Implémenter la structure de base d'une matrice en Python pur et réaliser une fonction de produit matriciel. Cela servira de base pour construire des matrices de passage.
\subsection*{Code Python}
\begin{verbatim}
def zeros(rows, cols):
    return [[0.0 for _ in range(cols)] for _ in range(rows)]

def eye(n):
    M = zeros(n, n)
    for i in range(n):
        M[i][i] = 1.0
    return M

def mat_mul(A, B):
    rows_A = len(A)
    cols_A = len(A[0])
    rows_B = len(B)
    cols_B = len(B[0])
    assert cols_A == rows_B, "Dimensions incompatibles pour le produit matriciel."

    C = zeros(rows_A, cols_B)
    for i in range(rows_A):
        for j in range(cols_B):
            s = 0.0
            for k in range(cols_A):
                s += A[i][k] * B[k][j]
            C[i][j] = s
    return C

def transpose(A):
    rows = len(A)
    cols = len(A[0])
    return [[A[i][j] for i in range(rows)] for j in range(cols)]

\section*{Tests de validation}
if __name__ == "__main__":
    A = [[1, 2], [3, 4]]
    I = eye(2)
    assert mat_mul(A, I) == A
    assert mat_mul(I, A) == A

    B = [[2, 0], [1, 2]]
    C = mat_mul(A, B)
    assert C == [[4, 4], [10, 8]]
    print("TP 01 validé.")
\end{verbatim}






\subsection*{Implémentation en Python Pur (Zéro Ellipse Algorithmique)}

\begin{verbatim}
def determinant_2x2(matrix):
    \"\"\"
    Calcule le determinant d'une matrice 2x2.
    Hypothese : matrix est une liste de 2 listes de taille 2 contenant des reels.
    \"\"\"
    return matrix[0][0] * matrix[1][1] - matrix[0][1] * matrix[1][0]

def invert_2x2(matrix):
    \"\"\"
    Inverse une matrice 2x2 en Python pur.
    \"\"\"
    det = determinant_2x2(matrix)
    if det == 0:
        raise ValueError("La matrice n'est pas inversible (determinant nul).")

    # Matrice des cofacteurs transposee (Comatrice)
    return [
        [matrix[1][1] / det, -matrix[0][1] / det],
        [-matrix[1][0] / det, matrix[0][0] / det]
    ]

def multiply_matrix_vector(matrix, vector):
    \"\"\"
    Multiplie une matrice nxn par un vecteur nx1.
    \"\"\"
    n = len(matrix)
    result = [0] * n
    for i in range(n):
        for j in range(n):
            result[i] += matrix[i][j] * vector[j]
    return result

\section*{--- Tests Mathematiques Rigoureux ---}
\section*{Matrice de passage de l'exercice :}
P = [[1, -1], [2, 1]]
\section*{Vecteur X dans B :}
X = [2, 5]

\section*{1. Inversion de P}
P_inv = invert_2x2(P)
\section*{Verification de l'inverse : P_inv devrait etre [[1/3, 1/3], [-2/3, 1/3]]}
assert abs(P_inv[0][0] - 1/3) < 1e-9
assert abs(P_inv[0][1] - 1/3) < 1e-9

\section*{2. Calcul des nouvelles coordonnees X' = P_inv * X}
X_prime = multiply_matrix_vector(P_inv, X)
\section*{X_prime devrait etre [7/3, 1/3]}
assert abs(X_prime[0] - 7/3) < 1e-9
assert abs(X_prime[1] - 1/3) < 1e-9

print("Toutes les assertions mathematiques sont verifiees avec succes.")
\end{verbatim}


\section*{TP 02}
\subsection*{Objectif}
Implémenter l'algorithme du pivot de Gauss pour inverser une matrice. C'est l'opération fondamentale pour passer de $P$ à $P^{-1}$ lors d'un changement de base.
\subsection*{Code Python}
\begin{verbatim}
def copy_mat(M):
    return [[val for val in row] for row in M]

def eye(n):
    return [[1.0 if i == j else 0.0 for j in range(n)] for i in range(n)]

def inverse(M):
    n = len(M)
    assert n == len(M[0]), "La matrice doit être carrée."
    A = copy_mat(M)
    Inv = eye(n)

    for i in range(n):
        # Recherche du pivot
        pivot_row = i
        for k in range(i + 1, n):
            if abs(A[k][i]) > abs(A[pivot_row][i]):
                pivot_row = k

        assert abs(A[pivot_row][i]) > 1e-9, "Matrice non inversible."

        # Échange des lignes
        A[i], A[pivot_row] = A[pivot_row], A[i]
        Inv[i], Inv[pivot_row] = Inv[pivot_row], Inv[i]

        # Normalisation du pivot
        pivot_val = A[i][i]
        for j in range(n):
            A[i][j] /= pivot_val
            Inv[i][j] /= pivot_val

        # Élimination
        for k in range(n):
            if k != i:
                factor = A[k][i]
                for j in range(n):
                    A[k][j] -= factor * A[i][j]
                    Inv[k][j] -= factor * Inv[i][j]
    return Inv

\section*{Fonction pour tester l'égalité avec tolérance}
def mat_approx_equal(A, B, tol=1e-6):
    for rA, rB in zip(A, B):
        for valA, valB in zip(rA, rB):
            if abs(valA - valB) > tol:
                return False
    return True

def mat_mul(A, B):
    n = len(A)
    return [[sum(A[i][k] * B[k][j] for k in range(n)) for j in range(n)] for i in range(n)]

\section*{Tests de validation}
if __name__ == "__main__":
    P = [[2, 3], [1, 2]]
    P_inv = inverse(P)

    I = mat_mul(P, P_inv)
    assert mat_approx_equal(I, eye(2)), "L'inverse calculé est incorrect."

    # Test avec l'exercice 3 : P = [[2,3],[1,2]], P_inv = [[2,-3],[-1,2]]
    assert mat_approx_equal(P_inv, [[2, -3], [-1, 2]])
    print("TP 02 validé.")
\end{verbatim}






\subsection*{Implémentation en Python Pur (Zéro Ellipse Algorithmique)}

\begin{verbatim}
def determinant_2x2(matrix):
    \"\"\"
    Calcule le determinant d'une matrice 2x2.
    Hypothese : matrix est une liste de 2 listes de taille 2 contenant des reels.
    \"\"\"
    return matrix[0][0] * matrix[1][1] - matrix[0][1] * matrix[1][0]

def invert_2x2(matrix):
    \"\"\"
    Inverse une matrice 2x2 en Python pur.
    \"\"\"
    det = determinant_2x2(matrix)
    if det == 0:
        raise ValueError("La matrice n'est pas inversible (determinant nul).")

    # Matrice des cofacteurs transposee (Comatrice)
    return [
        [matrix[1][1] / det, -matrix[0][1] / det],
        [-matrix[1][0] / det, matrix[0][0] / det]
    ]

def multiply_matrix_vector(matrix, vector):
    \"\"\"
    Multiplie une matrice nxn par un vecteur nx1.
    \"\"\"
    n = len(matrix)
    result = [0] * n
    for i in range(n):
        for j in range(n):
            result[i] += matrix[i][j] * vector[j]
    return result

\section*{--- Tests Mathematiques Rigoureux ---}
\section*{Matrice de passage de l'exercice :}
P = [[1, -1], [2, 1]]
\section*{Vecteur X dans B :}
X = [2, 5]

\section*{1. Inversion de P}
P_inv = invert_2x2(P)
\section*{Verification de l'inverse : P_inv devrait etre [[1/3, 1/3], [-2/3, 1/3]]}
assert abs(P_inv[0][0] - 1/3) < 1e-9
assert abs(P_inv[0][1] - 1/3) < 1e-9

\section*{2. Calcul des nouvelles coordonnees X' = P_inv * X}
X_prime = multiply_matrix_vector(P_inv, X)
\section*{X_prime devrait etre [7/3, 1/3]}
assert abs(X_prime[0] - 7/3) < 1e-9
assert abs(X_prime[1] - 1/3) < 1e-9

print("Toutes les assertions mathematiques sont verifiees avec succes.")
\end{verbatim}


\section*{TP 03}
\subsection*{Objectif}
Implémenter la formule mathématique du changement de base pour un vecteur : $X_{nouvelle} = P^{-1} X_{ancienne}$.
\subsection*{Code Python}
\begin{verbatim}
\section*{Fonctions des TP précédents}
def eye(n):
    return [[1.0 if i == j else 0.0 for j in range(n)] for i in range(n)]

def inverse(M):
    n = len(M)
    A = [[val for val in row] for row in M]
    Inv = eye(n)
    for i in range(n):
        pivot = i
        for k in range(i + 1, n):
            if abs(A[k][i]) > abs(A[pivot][i]): pivot = k
        A[i], A[pivot] = A[pivot], A[i]
        Inv[i], Inv[pivot] = Inv[pivot], Inv[i]
        p_val = A[i][i]
        for j in range(n):
            A[i][j] /= p_val
            Inv[i][j] /= p_val
        for k in range(n):
            if k != i:
                f = A[k][i]
                for j in range(n):
                    A[k][j] -= f * A[i][j]
                    Inv[k][j] -= f * Inv[i][j]
    return Inv

def mat_vec_mul(M, v):
    rows = len(M)
    cols = len(M[0])
    assert cols == len(v)
    return [sum(M[i][j] * v[j] for j in range(cols)) for i in range(rows)]

def change_basis_vector(X_old, P):
    """
    Calcule les coordonnées X_new d'un vecteur exprimé par X_old
    dans l'ancienne base, sachant que P est la matrice de passage.
    Formule : X_old = P * X_new  =>  X_new = P_inv * X_old
    """
    P_inv = inverse(P)
    return mat_vec_mul(P_inv, X_old)

\section*{Tests de validation}
if __name__ == "__main__":
    # Exercice 3 : P = [[2,3],[1,2]], X_old = [1, -1] -> X_new = [5, -3]
    P = [[2, 3], [1, 2]]
    X_old = [1, -1]
    X_new = change_basis_vector(X_old, P)

    # Tolérance pour la comparaison flottante
    assert abs(X_new[0] - 5.0) < 1e-6
    assert abs(X_new[1] - (-3.0)) < 1e-6

    # Exercice 1 (Cours) : P = [[1,1],[1,-1]], X_old = [4,2] -> X_new = [3,1]
    P2 = [[1, 1], [1, -1]]
    X_old2 = [4, 2]
    X_new2 = change_basis_vector(X_old2, P2)
    assert abs(X_new2[0] - 3.0) < 1e-6
    assert abs(X_new2[1] - 1.0) < 1e-6
    print("TP 03 validé.")
\end{verbatim}






\subsection*{Implémentation en Python Pur (Zéro Ellipse Algorithmique)}

\begin{verbatim}
def determinant_2x2(matrix):
    \"\"\"
    Calcule le determinant d'une matrice 2x2.
    Hypothese : matrix est une liste de 2 listes de taille 2 contenant des reels.
    \"\"\"
    return matrix[0][0] * matrix[1][1] - matrix[0][1] * matrix[1][0]

def invert_2x2(matrix):
    \"\"\"
    Inverse une matrice 2x2 en Python pur.
    \"\"\"
    det = determinant_2x2(matrix)
    if det == 0:
        raise ValueError("La matrice n'est pas inversible (determinant nul).")

    # Matrice des cofacteurs transposee (Comatrice)
    return [
        [matrix[1][1] / det, -matrix[0][1] / det],
        [-matrix[1][0] / det, matrix[0][0] / det]
    ]

def multiply_matrix_vector(matrix, vector):
    \"\"\"
    Multiplie une matrice nxn par un vecteur nx1.
    \"\"\"
    n = len(matrix)
    result = [0] * n
    for i in range(n):
        for j in range(n):
            result[i] += matrix[i][j] * vector[j]
    return result

\section*{--- Tests Mathematiques Rigoureux ---}
\section*{Matrice de passage de l'exercice :}
P = [[1, -1], [2, 1]]
\section*{Vecteur X dans B :}
X = [2, 5]

\section*{1. Inversion de P}
P_inv = invert_2x2(P)
\section*{Verification de l'inverse : P_inv devrait etre [[1/3, 1/3], [-2/3, 1/3]]}
assert abs(P_inv[0][0] - 1/3) < 1e-9
assert abs(P_inv[0][1] - 1/3) < 1e-9

\section*{2. Calcul des nouvelles coordonnees X' = P_inv * X}
X_prime = multiply_matrix_vector(P_inv, X)
\section*{X_prime devrait etre [7/3, 1/3]}
assert abs(X_prime[0] - 7/3) < 1e-9
assert abs(X_prime[1] - 1/3) < 1e-9

print("Toutes les assertions mathematiques sont verifiees avec succes.")
\end{verbatim}


\section*{TP 04}
\subsection*{Objectif}
Implémenter la formule mathématique de changement de base pour un endomorphisme : $A' = P^{-1} A P$.
\subsection*{Code Python}
\begin{verbatim}
def eye(n): return [[1.0 if i == j else 0.0 for j in range(n)] for i in range(n)]

def inverse(M):
    n = len(M)
    A = [[val for val in row] for row in M]
    Inv = eye(n)
    for i in range(n):
        pivot = i
        for k in range(i + 1, n):
            if abs(A[k][i]) > abs(A[pivot][i]): pivot = k
        A[i], A[pivot] = A[pivot], A[i]
        Inv[i], Inv[pivot] = Inv[pivot], Inv[i]
        p_val = A[i][i]
        for j in range(n):
            A[i][j] /= p_val; Inv[i][j] /= p_val
        for k in range(n):
            if k != i:
                f = A[k][i]
                for j in range(n):
                    A[k][j] -= f * A[i][j]; Inv[k][j] -= f * Inv[i][j]
    return Inv

def mat_mul(A, B):
    n = len(A)
    return [[sum(A[i][k] * B[k][j] for k in range(n)) for j in range(n)] for i in range(n)]

def change_basis_matrix(A_old, P):
    """
    Calcule la matrice de l'endomorphisme dans la nouvelle base.
    Formule : A_new = P^{-1} * A_old * P
    """
    P_inv = inverse(P)
    # A_new = P_inv * (A_old * P)
    AP = mat_mul(A_old, P)
    A_new = mat_mul(P_inv, AP)
    return A_new

def mat_approx_equal(A, B, tol=1e-6):
    for rA, rB in zip(A, B):
        for valA, valB in zip(rA, rB):
            if abs(valA - valB) > tol: return False
    return True

\section*{Tests de validation}
if __name__ == "__main__":
    # Test tiré de l'exercice 6
    A_old = [
        [2, -1, -1],
        [-1, 2, -1],
        [-1, -1, 2]
    ]
    P = [
        [1, 1, 1],
        [1, -1, 0],
        [1, 0, -1]
    ]
    A_new = change_basis_matrix(A_old, P)

    A_expected = [
        [0, 0, 0],
        [0, 3, 0],
        [0, 0, 3]
    ]
    assert mat_approx_equal(A_new, A_expected), "Le changement de base de l'endomorphisme a échoué."
    print("TP 04 validé.")
\end{verbatim}






\subsection*{Implémentation en Python Pur (Zéro Ellipse Algorithmique)}

\begin{verbatim}
def determinant_2x2(matrix):
    \"\"\"
    Calcule le determinant d'une matrice 2x2.
    Hypothese : matrix est une liste de 2 listes de taille 2 contenant des reels.
    \"\"\"
    return matrix[0][0] * matrix[1][1] - matrix[0][1] * matrix[1][0]

def invert_2x2(matrix):
    \"\"\"
    Inverse une matrice 2x2 en Python pur.
    \"\"\"
    det = determinant_2x2(matrix)
    if det == 0:
        raise ValueError("La matrice n'est pas inversible (determinant nul).")

    # Matrice des cofacteurs transposee (Comatrice)
    return [
        [matrix[1][1] / det, -matrix[0][1] / det],
        [-matrix[1][0] / det, matrix[0][0] / det]
    ]

def multiply_matrix_vector(matrix, vector):
    \"\"\"
    Multiplie une matrice nxn par un vecteur nx1.
    \"\"\"
    n = len(matrix)
    result = [0] * n
    for i in range(n):
        for j in range(n):
            result[i] += matrix[i][j] * vector[j]
    return result

\section*{--- Tests Mathematiques Rigoureux ---}
\section*{Matrice de passage de l'exercice :}
P = [[1, -1], [2, 1]]
\section*{Vecteur X dans B :}
X = [2, 5]

\section*{1. Inversion de P}
P_inv = invert_2x2(P)
\section*{Verification de l'inverse : P_inv devrait etre [[1/3, 1/3], [-2/3, 1/3]]}
assert abs(P_inv[0][0] - 1/3) < 1e-9
assert abs(P_inv[0][1] - 1/3) < 1e-9

\section*{2. Calcul des nouvelles coordonnees X' = P_inv * X}
X_prime = multiply_matrix_vector(P_inv, X)
\section*{X_prime devrait etre [7/3, 1/3]}
assert abs(X_prime[0] - 7/3) < 1e-9
assert abs(X_prime[1] - 1/3) < 1e-9

print("Toutes les assertions mathematiques sont verifiees avec succes.")
\end{verbatim}


\section*{TP 05}
\subsection*{Objectif}
Implémenter des opérations sur des matrices par blocs, spécifiquement le produit par blocs, sans utiliser le produit matriciel global (pour simuler un système distribué).
\subsection*{Code Python}
\begin{verbatim}
def zeros(r, c): return [[0.0]*c for _ in range(r)]

def mat_mul(A, B):
    rA, cA, cB = len(A), len(A[0]), len(B[0])
    C = zeros(rA, cB)
    for i in range(rA):
        for j in range(cB):
            C[i][j] = sum(A[i][k] * B[k][j] for k in range(cA))
    return C

def mat_add(A, B):
    return [[A[i][j] + B[i][j] for j in range(len(A[0]))] for i in range(len(A))]

def block_mul(A11, A12, A21, A22, B11, B12, B21, B22):
    """
    Calcule le produit de deux matrices 2x2 par blocs.
    C11 = A11*B11 + A12*B21
    C12 = A11*B12 + A12*B22
    C21 = A21*B11 + A22*B21
    C22 = A21*B12 + A22*B22
    """
    C11 = mat_add(mat_mul(A11, B11), mat_mul(A12, B21))
    C12 = mat_add(mat_mul(A11, B12), mat_mul(A12, B22))
    C21 = mat_add(mat_mul(A21, B11), mat_mul(A22, B21))
    C22 = mat_add(mat_mul(A21, B12), mat_mul(A22, B22))
    return C11, C12, C21, C22

def construct_from_blocks(C11, C12, C21, C22):
    n = len(C11)
    C = zeros(2*n, 2*n)
    for i in range(n):
        for j in range(n):
            C[i][j] = C11[i][j]
            C[i][j+n] = C12[i][j]
            C[i+n][j] = C21[i][j]
            C[i+n][j+n] = C22[i][j]
    return C

\section*{Tests de validation}
if __name__ == "__main__":
    I = [[1, 0], [0, 1]]
    Z = [[0, 0], [0, 0]]
    A12 = [[1, 2], [3, 4]]

    # M = [I  A12]
    #     [Z   I ]
    # M^2 = [I  2*A12]
    #       [Z   I   ]
    C11, C12, C21, C22 = block_mul(I, A12, Z, I, I, A12, Z, I)

    assert C11 == I
    assert C21 == Z
    assert C22 == I
    assert C12 == [[2, 4], [6, 8]]

    full_M2 = construct_from_blocks(C11, C12, C21, C22)
    assert full_M2 == [
        [1, 0, 2, 4],
        [0, 1, 6, 8],
        [0, 0, 1, 0],
        [0, 0, 0, 1]
    ]
    print("TP 05 validé.")
\end{verbatim}






\subsection*{Implémentation en Python Pur (Zéro Ellipse Algorithmique)}

\begin{verbatim}
def determinant_2x2(matrix):
    \"\"\"
    Calcule le determinant d'une matrice 2x2.
    Hypothese : matrix est une liste de 2 listes de taille 2 contenant des reels.
    \"\"\"
    return matrix[0][0] * matrix[1][1] - matrix[0][1] * matrix[1][0]

def invert_2x2(matrix):
    \"\"\"
    Inverse une matrice 2x2 en Python pur.
    \"\"\"
    det = determinant_2x2(matrix)
    if det == 0:
        raise ValueError("La matrice n'est pas inversible (determinant nul).")

    # Matrice des cofacteurs transposee (Comatrice)
    return [
        [matrix[1][1] / det, -matrix[0][1] / det],
        [-matrix[1][0] / det, matrix[0][0] / det]
    ]

def multiply_matrix_vector(matrix, vector):
    \"\"\"
    Multiplie une matrice nxn par un vecteur nx1.
    \"\"\"
    n = len(matrix)
    result = [0] * n
    for i in range(n):
        for j in range(n):
            result[i] += matrix[i][j] * vector[j]
    return result

\section*{--- Tests Mathematiques Rigoureux ---}
\section*{Matrice de passage de l'exercice :}
P = [[1, -1], [2, 1]]
\section*{Vecteur X dans B :}
X = [2, 5]

\section*{1. Inversion de P}
P_inv = invert_2x2(P)
\section*{Verification de l'inverse : P_inv devrait etre [[1/3, 1/3], [-2/3, 1/3]]}
assert abs(P_inv[0][0] - 1/3) < 1e-9
assert abs(P_inv[0][1] - 1/3) < 1e-9

\section*{2. Calcul des nouvelles coordonnees X' = P_inv * X}
X_prime = multiply_matrix_vector(P_inv, X)
\section*{X_prime devrait etre [7/3, 1/3]}
assert abs(X_prime[0] - 7/3) < 1e-9
assert abs(X_prime[1] - 1/3) < 1e-9

print("Toutes les assertions mathematiques sont verifiees avec succes.")
\end{verbatim}


\end{document}
